% !TeX root = ../../Book2.tex
\chapter{Bochner–Riesz 平均}
\label{b2:ch:24}
\BookChapterNumber{b2:ch:24}
\begin{chapterabstract}
Bochner--Riesz 平均通过球边界的消失阶调整频率截断。本章从符号缩放、阶数平均与核的振荡渐近出发，证明可积核范围、尾部必要条件、解析插值充分范围以及带正则性的逼近估计。球面限制提供另一类混合指数控制，Fefferman 定理则揭示高维硬球截断的几何障碍。半径、平滑阶数与输入空间由此呈现出不同的作用。
\end{chapterabstract}

\begin{learninggoals}\begin{itemize}
\item\label{b2:item:24-goal-1}核对半径缩放、阶数平均、Bessel 核常数及两参数极限，分清频率边界平滑与低频逼近阶数。
\item\label{b2:item:24-goal-2}用振荡主项和可控余项判断核的可积性与必要条件，理解端点为何需要下界而不只需要衰减上界。
\item\label{b2:item:24-goal-3}区分可积核范围、已证插值范围、经典目标与深层否定定理，并核对一致有界性在范数收敛中的作用。
\item\label{b2:item:24-goal-4}将限制估计用于径向壳层，保留混合指数和尺度因子；在组合平均或几何逼近中检查常数是否一致。
\end{itemize}\end{learninggoals}

矩形硬截断可以按坐标分解，球形硬截断却同时包含连续变化的法向。给球边界增加消失因子，能够减弱频率截断造成的远场振荡，但需要多少平滑仍取决于所问的空间。平方范数只看符号大小，端点空间要求核有足够的绝对可积性，中间指数则需要更细的抵消与几何估计。

\ProseParagraphBreak

学习本章时，可先分别观察两个参数。半径增大使更多频率进入平均，阶数增大则改变球内部的权重；两者同时变化甚至可以产生热算子或零算子，而不一定逼近恒等。\secref{b2:sec:2401}确定这些算子的定义与缩放，\secref{b2:sec:2402}再说明边界消失阶如何转化为空间核的衰减。

\ProseParagraphBreak

\secref{b2:sec:2403}将必要条件与充分估计放在同一组参数中比较。一个没有被尾部反例排除的指标，并不因此获得有界性；反过来，可积核方法或简单插值无法覆盖它，也不能据此宣布无界。逼近速度还需单独检查符号在原点的 Taylor 展开。组合不同半径可以消去低阶误差，却不会自动消除各条球边界的乘子障碍。

\ProseParagraphBreak

\appref{b2:app:A}的 Bessel 公式将给出径向核的具体表达，核的衰减又决定哪些指数可以直接使用 Young 不等式。最后一节将球内分解成许多薄层，把问题与\chapref{b2:ch:23}的限制估计相联系。Fefferman 定理说明球形硬截断的特殊障碍；其深层几何证明在此仅作概要介绍。

% !TeX root = ../../Book2.tex
\section{Bochner–Riesz 乘子}
\label{b2:sec:2401}

球形频率截断保留低频、删除高频，是一种自然的 Fourier 求和方法。它的困难集中在球面边界：边界两侧的符号突然跳变，逆变换便产生长距离振荡。Bochner–Riesz 平均在球内加入一个向边界逐渐消失的权，用参数 \(\delta\) 记录消失阶数。需要确定的是，对给定的 \(L^p\) 空间，多少平滑足以同时保证算子有界和半径增大时的逼近。

\subsection{定义}
\label{b2:subsec:240101}

本章采用
\(\widehat f(\xi)=\int_{\mathbb R^d}f(x)e^{-2\pi\mathrm i x\cdot\xi}\,dx\)。
先在有界符号能够直接作用的平方空间上定义算子，再讨论其他指数。

\Needspace{6\baselineskip}
\begin{definition}[Bochner--Riesz 平均]\label{b2:def:24-bochner-riesz}
设整数 \(d\geq1\)、\(\delta\geq0\)、\(R>0\)。定义频率符号
\[
 m_{\delta,R}(\xi)\coloneqq
 \begin{cases}
 (1-|\xi|^2/R^2)^\delta,&|\xi|<R,\\
 0,&|\xi|\geq R.
 \end{cases}
\]
相应的\term[bochner-riesz-pingjun]{Bochner--Riesz 平均}[Bochner--Riesz means]是
\[
 B_R^\delta f\coloneqq
       \mathcal F^{-1}(m_{\delta,R}\cdot\widehat f),
 \quad f\in L^2(\mathbb R^d).
\]
单位半径核记为
\[
 K_\delta\coloneqq\mathcal F^{-1}m_{\delta,1}.
\]
这里 \(\delta=0\) 时的符号是开球示性函数；球面上的取值对这个
Lebesgue 乘子没有影响。
\end{definition}
\symidx{\(B_R^\delta f\)：Bochner--Riesz 平均}
\symidx{\(K_\delta\)：单位半径 Bochner--Riesz 核}

由于 \(0\leq m_{\delta,R}\leq1\)，该定义在 \(L^2\) 中总有意义。
对 Schwartz 输入，还有普通积分表达式
\[
 B_R^\delta f(x)
 =\int_{|\xi|<R}e^{2\pi\mathrm i x\cdot\xi}
           (1-|\xi|^2/R^2)^\delta\widehat f(\xi)\,d\xi.
\]
它处处绝对收敛，输出光滑。若要在 \(1\leq p<\infty\) 上定义同一算子，
则须先证明
\[
 \|B_R^\delta f\|_p\leq C\|f\|_p
       \quad(f\in\mathcal S).
\]
Schwartz 类的稠密性随后给出唯一连续延拓，并在 \(L^p\cap L^2\)
上与平方定义一致。无穷指数采用\chapref{b2:ch:22}的伴随约定：
从相同实符号的 \(L^1\) 算子取 Banach 伴随；核可积时，这一约定
就是对有界输入作卷积。不能改用 Schwartz 类在 \(L^\infty\) 中的范数稠密性。

\ProseParagraphBreak
正参数使符号在球面处连续，但不使它在边界任意光滑。
非整数 \(\delta\) 的法向导数在超过相应阶数后出现奇性；
正整数阶也会在更高导数处遇到跨边界的不匹配。
因此，不能把“边界值连续”直接替换成\chapref{b2:ch:22} Mihlin 判据所需的全部导数控制。
稍后将用实际核估计确定有界范围。

\ProseParagraphBreak
这个算子整合球内部的频率，并不是只在球面上积分的延拓算子。
两者的联系来自逐半径分层。先取
\(\chi\in C_c^\infty(\mathbb R^d)\)，使 \(\chi=1\) 于原点附近，且支撑
严格位于单位球内。符号 \(\chi\cdot m_{\delta,1}\) 光滑紧支撑，其逆变换
为 Schwartz 函数，故相应卷积在所有 \(L^p\) 上有界。剩余部分才包含
球面边界的困难。这一区分会在核渐近、解析插值和薄壳分解中反复出现。

\subsection{缩放}
\label{b2:subsec:240102}

半径变化首先是一个准确的共轭关系。为了不混淆正、反缩放，明确规定
\[
 D_Rf(x)\coloneqq f(Rx),\quad R>0.
\]
它满足
\(\widehat{D_Rf}(\xi)=R^{-d}\widehat f(\xi/R)\) 和
\(\|D_Rf\|_p=R^{-d/p}\|f\|_p\)，无穷指数按 \(1/p=0\) 理解。

\Needspace{6\baselineskip}
\begin{proposition}[半径共轭与一致常数]\label{b2:prop:24-dilation}
设 \(d\geq1\)、\(\delta\geq0\)，\(B_R^\delta\)、\(K_\delta\)
如\cref{b2:def:24-bochner-riesz}。则对 Schwartz 函数有
\[
 B_R^\delta=D_R B_1^\delta D_{1/R},\quad
 K_{\delta,R}(x)=R^dK_\delta(Rx).
\]
若 \(B_1^\delta\) 在 \(L^p(\mathbb R^d)\)、\(1\leq p<\infty\)
上有界，则所有 \(R>0\) 的算子都有唯一相容延拓，且
\[
 \|B_R^\delta\|_{L^p\to L^p}
       =\|B_1^\delta\|_{L^p\to L^p}.
\]
在 \(L^\infty\) 上使用 \(L^1\) 伴随延拓时，同一等式成立。
\end{proposition}
\begin{proof}
先对 Schwartz 输入作 Fourier 变换：
\[
 \begin{aligned}
 \mathcal F(D_R B_1^\delta D_{1/R}f)(\xi)
 &=R^{-d}m_{\delta,1}(\xi/R)\,
             R^d\widehat f(\xi)\\
 &=m_{\delta,R}(\xi)\widehat f(\xi).
 \end{aligned}
\]
核公式则直接由逆变换中的 \(\xi=R\eta\) 得到。
两个拉回的范数因子互为倒数，所以共轭给出
\(\|B_R^\delta\|\leq\|B_1^\delta\|\)。反向用
\(B_1^\delta=D_{1/R}B_R^\delta D_R\)，得到反向不等式。
有限指数上的稠密延拓保留恒等式。无穷指数从 \(L^1\) 的共轭恒等式
取伴随，或在核可积时直接换元，得到相同结论。
\end{proof}

所以固定 \(\delta,p\) 的问题不是分别控制每个半径，而是证明一个
单位半径界。这个界一旦成立，就给出整个族的统一常数；缩放本身不能
证明单位半径有界。

\Needspace{6\baselineskip}
\begin{proposition}[有界性推出范数逼近]\label{b2:prop:24-strong-approximation}
设 \(d\geq1\)、\(\delta\geq0\)、\(1\leq p<\infty\)，并假设
\cref{b2:def:24-bochner-riesz}的 \(B_1^\delta\) 有界于
\(L^p(\mathbb R^d)\)。则对每个 \(f\in L^p(\mathbb R^d)\)，
\[
 \lim_{R\to\infty}\|B_R^\delta f-f\|_p=0.
\]
\end{proposition}
\begin{proof}
先取 \(\widehat h\in C_c^\infty\)，支撑包含在固定球中。
当 \(R\) 充分大时，整个支撑远离 \(|\xi|=R\)，在这片固定频率上，
\[
 (m_{\delta,R}-1)\widehat h\longrightarrow0
       \quad\text{于 }C_c^\infty
\]
的所有导数半范数中。事实上，正阶时对
\((1-|\xi|^2/R^2)^\delta\) 求导，分母始终远离零，每个非恒定项
含有趋零的 \(R^{-2}\) 因子；零阶时误差在这个支撑上最终恒零。
固定紧支撑上的导数收敛经 Fourier 反演和分部积分给出
\(B_R^\delta h-h\to0\) 于 Schwartz 拓扑，特别在 \(L^p\) 中收敛。

这样的 \(h\) 在有限指数 \(L^p\) 中稠密。为看清这一点，先以
\(C_c^\infty\) 函数逼近任意输入，再对这个 Schwartz 函数 \(u\)
使用频率截断
\(\widehat h_N=\chi(\xi/N)\widehat u(\xi)\)，其中固定光滑 \(\chi\)
紧支撑并在原点附近等于一。Schwartz 衰减控制截去的频率尾部及各阶导数，
故 \(h_N\to u\) 于 Schwartz 拓扑。

记统一算子界为 \(C\)。给定 \(\varepsilon>0\)，选上述 \(h\)
使 \(\|f-h\|_p<\varepsilon\)。于是
\[
 \|B_R^\delta f-f\|_p
 \leq (C+1)\varepsilon+\|B_R^\delta h-h\|_p.
\]
先令 \(R\to\infty\)，再令 \(\varepsilon\downarrow0\)，得到结论。
\end{proof}

证明中真正把稠密类推广到全部输入的是统一算子界。仅知道光滑输入
最终逼近良好，并不能反推普通 \(L^p\) 的有界性。
另一方面，这是逐输入的范数收敛，也不是算子范数收敛：
在 \(L^2\) 中，\(I-B_R^\delta\) 的符号在球外恒为一、处处模不超过一，
所以
\[
 \|I-B_R^\delta\|_{2\to2}=1\quad\text{对每个 }R>0.
\]
把单位输入的频率放到当前球外，总能留下全部误差。
固定输入的频率尾部会消失，随半径移动的输入却没有这种统一性。

\ProseParagraphBreak
阶数方向则有另一种结构：较高阶可以写成较低阶按半径的正权平均。
记 beta 积分
\[
 \mathrm B(a,b)=\int_0^1t^{a-1}(1-t)^{b-1}\,dt
       =\frac{\Gamma(a)\Gamma(b)}{\Gamma(a+b)},\quad a,b>0.
\]

\Needspace{6\baselineskip}
\begin{proposition}[阶数的单调传递]\label{b2:prop:24-order-monotonicity}
设 \(d\geq1\)、\(0\leq\delta_1<\delta_2\)，置
\(\eta=\delta_2-\delta_1\)。对\cref{b2:def:24-bochner-riesz}
的平均和任意 Schwartz 输入，有
\[
 B_R^{\delta_2}f=
 \frac1{\mathrm B(\delta_1+1,\eta)}
 \int_0^1 B_{R\sqrt t}^{\delta_1}f\,
          t^{\delta_1}(1-t)^{\eta-1}\,dt,\quad R>0.
\]
若 \(1\leq p<\infty\) 且
\(\sup_{R>0}\|B_R^{\delta_1}\|_{p\to p}\leq C\)，则
\(\sup_{R>0}\|B_R^{\delta_2}\|_{p\to p}\leq C\)。
在共同的 Schwartz 测试类上还有点态不等式
\[
 \sup_{R>0}|B_R^{\delta_2}f(x)|
       \leq\sup_{R>0}|B_R^{\delta_1}f(x)|.
\]
\end{proposition}
\begin{proof}
令 \(a=|\xi|^2/R^2\)。在 \(a<1\) 时，右侧的符号为
\[
 \begin{aligned}
 &\frac1{\mathrm B(\delta_1+1,\eta)}
     \int_a^1(1-a/t)^{\delta_1}t^{\delta_1}
                         (1-t)^{\eta-1}\,dt\\
 &\quad=\frac1{\mathrm B(\delta_1+1,\eta)}
     \int_a^1(t-a)^{\delta_1}(1-t)^{\eta-1}\,dt
       =(1-a)^{\delta_2}.
 \end{aligned}
\]
最后用 \(t=a+(1-a)u\) 换元。若 \(a\geq1\)，两侧均为零。
对 Schwartz 输入，可以将半径积分移过逆 Fourier 变换：
符号的模至多为一，\(\widehat f\) 可积，而归一化 beta 权的总质量为一。
这证明积分恒等式及其点态意义。

若较低阶的 \(L^p\) 界已知，\(t\mapsto B_{R\sqrt t}^{\delta_1}f\)
在 \(t>0\) 强连续。由缩放恒等式，这归结为有限指数 \(L^p\)
上拉回 \(D_r\) 的强连续性，后者先对光滑紧支撑函数验证，再用稠密性。
因此加权积分也可在 \(L^p\) 中解释，且其范数由
\(C\|f\|_p\) 控制。Minkowski 与测试函数的稠密延拓给出较高阶的界。
对点态式取绝对值，再用总质量为一的正权积分，得到最后的不等式。
\end{proof}

正权只出现在半径平均中，不说明实空间核处处非负。
例如低阶核本来有正负振荡，按半径作正权平均以后仍可能变号。
这个公式只能把有界性向更高阶传递，不能从大阶数的可积核反推出
小阶数的结论。若要把极大不等式从测试类推广到一般输入，也仍须
使用已建立的极大算子延拓，不能任意选择各半径的代表后取上确界。

\ProseParagraphBreak
固定阶数时 \(R\to\infty\) 逼近恒等；阶数同时变化，则可能得到热流。
这个例子提醒我们，两个参数不能只观察其中一个。

\Needspace{6\baselineskip}
\begin{proposition}[阶数与半径的联合平方极限]\label{b2:prop:24-heat-limit}
设 \(d\geq1\)、\(R_n\to\infty\)、\(\delta_n\geq0\)，并假设
\(\delta_n/R_n^2\to a\in[0,\infty]\)。对每个
\(f\in L^2(\mathbb R^d)\)，\cref{b2:def:24-bochner-riesz}
中的平均满足
\[
 B_{R_n}^{\delta_n}f\longrightarrow
 \begin{cases}
 f,&a=0,\\
 e^{\,a\Delta/(4\pi^2)}f,&0<a<\infty,\\
 0,&a=\infty
 \end{cases}
 \quad\text{于 }L^2.
\]
中间的算子按 Fourier 符号 \(e^{-a|\xi|^2}\) 定义。
\end{proposition}
\begin{proof}
固定 \(\xi\)，当 \(n\) 充分大时 \(|\xi|<R_n\)。
若 \(a<\infty\)，则
\[
 \delta_n\log(1-|\xi|^2/R_n^2)
 =-\frac{\delta_n}{R_n^2}|\xi|^2
       +O\!\left(\frac{\delta_n}{R_n^4}|\xi|^4\right)
       \longrightarrow-a|\xi|^2.
\]
误差趋零，因为 \(\delta_n/R_n^2\) 有界而 \(R_n^{-2}\to0\)。
若 \(a=\infty\) 且 \(\xi\neq0\)，则
\[
 0\leq m_{\delta_n,R_n}(\xi)
       \leq\exp(-\delta_n|\xi|^2/R_n^2)\longrightarrow0.
\]
零频单点不影响 Lebesgue 积分。所有符号均介于零与一之间，
故平方误差由 \(4|\widehat f|^2\) 控制。Plancherel 和支配收敛
给出所述三种强极限。热算子时间是 \(a/(4\pi^2)\)，来自本书
Fourier 变换中 Laplace 算子的符号 \(-4\pi^2|\xi|^2\)。
\end{proof}

特别地，\(R_n=\sqrt{\delta_n}\to\infty\) 产生符号 \(e^{-|\xi|^2}\)。
这是 \(L^2\) 中的强算子极限；有限阶核是否非负、其他
\(L^p\) 上是否存在统一算子界，还取决于近似核本身的性质。

\subsection{核表示}
\label{b2:subsec:240103}

为了读出边界平滑的作用，现在计算单位半径核。
核积分可以比普通平均的参数范围更广：只要 \(\delta>-1\)，
球内幂权仍可积；但负参数的符号在边界无界，不能当作平方有界乘子。

\Needspace{6\baselineskip}
\begin{proposition}[Bochner--Riesz 核]\label{b2:prop:24-kernel}
设 \(d\geq1\)、\(\delta>-1\)，在开单位球内令
\(m_\delta(\xi)=(1-|\xi|^2)^\delta\)，球外取零。
零阶按球示性函数解释。则
\[
 K_\delta(x)\coloneqq\int_{|\xi|<1}
             e^{2\pi\mathrm i x\cdot\xi}m_\delta(\xi)\,d\xi
\]
是连续有界的实值径向函数，对 \(x\neq0\) 满足
\[
 K_\delta(x)=\Gamma(\delta+1)\pi^{-\delta}
       |x|^{-d/2-\delta}J_{d/2+\delta}(2\pi|x|).
\]
在原点，
\[
 K_\delta(0)=
       \frac{\pi^{d/2}\Gamma(\delta+1)}
            {\Gamma(\delta+d/2+1)}.
\]
\end{proposition}
\begin{proof}
极坐标下，球内符号在边界的可积性恰由
\(\int_{1/2}^1(1-r)^\delta\,dr<\infty\) 保证。
因此逆变换连续有界；旋转不变性给出径向性，频率偶对称使正弦部分抵消，
给出实值性。

Bessel 公式使用\cref{b2:thm:A2-bochner-riesz}，其所需条件正是
\(d\geq1,\delta>-1\)。回忆该处计算的关键一步：将径向变换中的
Bessel 级数逐项积分后，出现
\[
 \int_0^1(1-s^2)^\delta s^{2k+2\nu+1}\,ds
 =\frac{\Gamma(k+\nu+1)\Gamma(\delta+1)}
        {2\Gamma(k+\nu+\delta+2)},\quad \nu=d/2-1.
\]
它消去原级数中的一个 Gamma 分母，剩下阶数
\(\nu+\delta+1=d/2+\delta\) 的 Bessel 级数。固定空间半径时，
绝对值级数由可积球内权控制，故这次逐项积分合法。
代回 Fourier 归一化中的 \(2\pi\)，得到所写系数。

原点值也可不借 Bessel 极限直接检查：
\[
 \begin{aligned}
 K_\delta(0)
 &=\frac{2\pi^{d/2}}{\Gamma(d/2)}
             \int_0^1(1-r^2)^\delta r^{d-1}\,dr\\
 &=\frac{\pi^{d/2}}{\Gamma(d/2)}
             \mathrm B(d/2,\delta+1)
 =\frac{\pi^{d/2}\Gamma(\delta+1)}
             {\Gamma(\delta+d/2+1)}.
 \end{aligned}
\]
\end{proof}

例如 \(d=2,\delta=0\) 时
\(K_0(x)=J_1(2\pi|x|)/|x|\)，原点值为 \(\pi\)，正是单位圆盘面积。
核的 Bessel 阶数与球面测度不同：球面公式出现 \(d/2-1\)，
球内加权公式出现 \(d/2+\delta\)。球内部的径向积分及其边界消失阶
正是这一差别的来源。下一节把公式化为可用于绝对积分和反例的远场渐近。

% !TeX root = ../../Book2.tex
\section{核的渐近}
\label{b2:sec:2402}

核公式给出了一个特殊函数表达式，要判断可积性，还必须知道无穷远处的主项及误差。上界能够证明收敛，排除端点则需要下界。本节把球内积分化为两个一维端点的贡献，读出边界消失阶怎样增加空间衰减，再由同一渐近分别处理核范数、逼近恒等和极大控制。

\subsection{Bessel 表示}
\label{b2:subsec:240201}

除 Bessel 级数外，还有一个便于看清几何的切片表示。
设 \(d\geq2,\delta>-1\)，把空间点旋转为 \(x=(0,\ldots,0,r)\)。
频率写为 \((\eta,t)\in\mathbb R^{d-1}\times\mathbb R\)，则
\[
 \begin{aligned}
 K_\delta(r)
 &=\int_{-1}^1e^{2\pi\mathrm i rt}
       \int_{|\eta|^2<1-t^2}
          (1-t^2-|\eta|^2)^\delta\,d\eta\,dt.
 \end{aligned}
\]
所有换序都由球内幂权的绝对可积性保证。
内层令 \(\eta=(1-t^2)^{1/2}u\)，得到
\[
 \int_{|\eta|^2<1-t^2}(1-t^2-|\eta|^2)^\delta\,d\eta
 =c_{d,\delta}(1-t^2)^{\delta+(d-1)/2},
\]
其中用\cref{b2:prop:24-kernel}的原点积分在维数 \(d-1\) 中计算，
\[
 c_{d,\delta}
 =\frac{\pi^{(d-1)/2}\Gamma(\delta+1)}
        {\Gamma(\delta+(d+1)/2)}.
\]
故令 \(\alpha=\delta+(d-1)/2\)，有
\[
 K_\delta(r)=c_{d,\delta}
       \int_{-1}^1e^{2\pi\mathrm i rt}(1-t^2)^\alpha\,dt.
\]
当 \(d=1\) 时，直接从定义得到同一公式，此时 \(c_{1,\delta}=1\)。

\ProseParagraphBreak
横截面的尺度贡献 \((1-t^2)^{(d-1)/2}\)，原来的边界权再贡献
\((1-t^2)^\delta\)。两者相加，给出端点 \(t=\pm1\) 的消失阶
\(\alpha\)。区间内部的振幅光滑，且相位导数不为零，局部化后可
反复分部积分；真正决定主项的是两个端点。

\ProseParagraphBreak
低维模型可以同时检查常数与符号。先取 \(d=1\)，置 \(u=2\pi x\)。
零阶时直接积分得到
\[
 K_0(x)=\int_{-1}^1e^{\mathrm i u\xi}\,d\xi
       =\frac{2\sin u}{u},\quad K_0(0)=2.
\]
对一阶权 \(1-\xi^2\)，在有限区间上对 \(u\) 求导合法，因而
\[
 K_1(x)=\left(1+\frac{d^2}{du^2}\right)
              \frac{2\sin u}{u}
       =\frac{4(\sin u-u\cos u)}{u^3},\quad K_1(0)=\frac43.
\]
后一个原点值可由 Taylor 展开取得，也等于
\(\int_{-1}^1(1-\xi^2)\,d\xi\)。
当 \(u=2\pi\) 时，\(K_1(x)<0\)，所以即使正阶平均的核已可积，
它也不是一个非负的概率核。

\ProseParagraphBreak
在三维，零阶球的横截面积为 \(\pi(1-t^2)\)，于是
\[
 K_0(x)=4\pi\,\frac{\sin u-u\cos u}{u^3},
 \quad u=2\pi|x|,\quad K_0(0)=\frac{4\pi}{3}.
\]
它与一维一阶核只差常数，但绝对积分的径向体积元已经由 \(dr\)
变为 \(4\pi r^2\,dr\)。同一个函数形式在不同底空间有不同的
可积性，判断时不能省去体积增长。

\subsection{端点振荡渐近}
\label{b2:subsec:240202}

所需的一维工具是\cref{b2:lem:A2-endpoint-asymptotic}：
若 \(\beta>-1\)，\(\chi\) 在 \([0,\infty)\) 光滑且紧支撑，则
\[
 \int_0^\infty e^{-2\pi\mathrm i rs}s^\beta\chi(s)\,ds
 =\Gamma(\beta+1)(2\pi r)^{-\beta-1}
       e^{-\mathrm i\pi(\beta+1)/2}\chi(0)
       +O(r^{-\beta-2}).
\]
这里 \(r\to+\infty\)，复相位对应负号 Fourier 指数。
该结论不要求 \(\beta\) 为整数，因而能处理任意实消失阶。

\ProseParagraphBreak
回忆这个工具的证明机制，可以避免把端点公式误当成一次普通分部积分。
模型积分带上 \(e^{-s}\) 后为
\[
 \int_0^\infty e^{-2\pi\mathrm i rs}s^\beta e^{-s}\,ds
       =\Gamma(\beta+1)(1+2\pi\mathrm i r)^{-\beta-1}.
\]
将 \(\chi(s)\) 写成
\(\chi(0)e^{-s}+s\,c(s)\)，余项在零点多消失一阶。
在 \(0\leq s\lesssim r^{-1}\) 上直接积分，得到 \(O(r^{-\beta-2})\)；
在 \(s\asymp2^j/r\) 的剩余小环带上，取整数 \(N>\beta+2\)
次分部积分，得到
\[
 C r^{-\beta-2}2^{j(\beta+2-N)}.
\]
这是可求和的几何级数，远离零点的部分则有任意阶衰减。
这个分解既利用了端点的尺度，又没有对奇异振幅在零点作非法求导。
附录中的证明同时控制了截断导数与无穷远的指数尾部。

\ProseParagraphBreak
现在固定 \(d\geq1,\delta\geq0\)，将切片积分在两个端点与区间内部
作光滑分割。在 \(t=1\) 附近令 \(s=1-t\)，便有
\[
 (1-t^2)^\alpha=s^\alpha(2-s)^\alpha,\quad
 e^{2\pi\mathrm i rt}=e^{2\pi\mathrm i r}e^{-2\pi\mathrm i rs}.
\]
端点振幅在 \(s=0\) 的值为 \(2^\alpha\)。上述公式给出的正端点主项为
\[
 c_{d,\delta}\,2^\alpha\Gamma(\alpha+1)
       (2\pi r)^{-\alpha-1}
       e^{\,2\pi\mathrm i r-\mathrm i\pi(\alpha+1)/2}.
\]
负端点贡献是它的共轭，内部项经分部积分可忽略到所需阶数。
代入 \(\alpha=\delta+(d-1)/2\)，整理常数得到
\[
 \begin{aligned}
 K_\delta(r)
 &=a_\delta r^{-\delta-(d+1)/2}
       \cos\!\left(2\pi r-\frac{\pi(d+1+2\delta)}4\right)
       +O(r^{-\delta-(d+3)/2}),\\
 a_\delta&=\frac{\Gamma(\delta+1)}{\pi^{\delta+1}}>0.
 \end{aligned}
\]
这也与\cref{b2:prop:24-kernel}中 Bessel 阶数
\(\nu=d/2+\delta\) 的大半径渐近一致。

\ProseParagraphBreak
每增加一阶 \(\delta\)，空间主项便多衰减一个幂次。
在固定维数下，若 \(\delta\) 限制在 \([0,\infty)\) 的一个固定紧区间，
这里的余项常数还可以一致选取：端点分割固定，光滑因子
\((2-s)^\alpha\) 的有限阶导数一致有界；环带估计中的 \(N\) 取为
大于全部 \(\alpha+2\) 的同一个整数，幂函数积分的分母也远离零。
Gamma 系数在该参数紧集连续有界。这种局部参数一致性可用于研究
可积性阈值附近的范数变化，却不是允许阶数无界增长的渐近结论。

\ProseParagraphBreak
余弦零点并不会使主项在大部分半径上消失。令
\[
 \beta_\delta=\delta+\frac{d+1}{2},\quad
 \phi_\delta=\frac{\pi(d+1+2\delta)}4.
\]
在每个充分大的单位区间 \([n,n+1]\) 内，集合
\[
 A_n=\{r\in[n,n+1]:|\cos(2\pi r-\phi_\delta)|\geq3/4\}
\]
都有同一个正长度。余项比主振幅小一个 \(r^{-1}\) 因子，
因此当 \(n\) 足够大时，
\[
 |K_\delta(r)|\geq c_{d,\delta}r^{-\beta_\delta}
       \quad(r\in A_n).
\]
上界处理整个空间，下界只需占固定比例的这些径向区域。
下面将这两种用途分开应用。

\subsection{空间衰减}
\label{b2:subsec:240203}

\Needspace{6\baselineskip}
\begin{proposition}[核的有限指数可积性]\label{b2:prop:24-kernel-tail}
设 \(d\geq1,\delta\geq0\)，\(K_\delta\) 为
\cref{b2:def:24-bochner-riesz}的单位半径核。
对每个 \(1\leq q<\infty\)，有
\[
 K_\delta\in L^q(\mathbb R^d)
 \quad\Longleftrightarrow\quad
 q\left(\delta+\frac{d+1}{2}\right)>d.
\]
核在全空间有界；上述有限指数条件在等号时不成立。
\end{proposition}
\begin{proof}
核在有界集合上连续有界。远场上界给出
\[
 \int_{|x|>1}|K_\delta(x)|^q\,dx
 \leq C\int_1^\infty r^{d-1-q\beta_\delta}\,dr,
 \quad \beta_\delta=\delta+(d+1)/2.
\]
当 \(q\beta_\delta>d\) 时收敛。
反过来，在前述 \(A_n\) 上使用下界及极坐标，
\[
 \int_{|x|>1}|K_\delta(x)|^q\,dx
       \geq c\sum_{n\geq n_0}n^{d-1-q\beta_\delta}.
\]
这里 \(r\asymp n\)，且每个 \(A_n\) 的长度都有固定正下界。
若 \(q\beta_\delta\leq d\)，这个级数发散；临界等号时就是调和级数。
\end{proof}

这只判断一个核函数的 \(L^q\) 性质。乘子在 \(L^p\) 上有界是对
全部输入的要求，两者不能互换。核的 \(L^1\) 可积性能够通过 Young
提供所有指数的充分条件，但不是每个中间指数的必要条件。

\Needspace{6\baselineskip}
\begin{theorem}[可积核范围]\label{b2:thm:24-integrable}
设 \(d\geq1,\delta\geq0\)，\(B_R^\delta,K_\delta\)
如\cref{b2:def:24-bochner-riesz}。则
\[
 K_\delta\in L^1(\mathbb R^d)
       \quad\Longleftrightarrow\quad \delta>\frac{d-1}{2}.
\]
在这一范围，算子由 \(K_{\delta,R}*f\) 定义于所有
\(L^p(\mathbb R^d)\)、\(1\leq p\leq\infty\)，并有
\[
 \sup_{R>0}\|B_R^\delta f\|_p
       \leq\|K_\delta\|_1\cdot\|f\|_p.
\]
对 \(p<\infty\)，\(B_R^\delta f\to f\) 于 \(L^p\)；
对有界一致连续函数，收敛在全空间一致范数中成立。
\end{theorem}
\begin{proof}
在上一命题中取 \(q=1\)，得到准确的核可积阈值。
缩放不改变核的 \(L^1\) 范数，Young 不等式遂给出算子界。
这个卷积与 Schwartz 类上的 Fourier 定义一致，有限指数上的延拓
因而相容；无穷指数也与 \(L^1\) 伴随约定一致，因为核实值且为偶函数。

还须核对逼近恒等的总质量。符号与核均可积，Fourier 反演表明
\(\widehat K_\delta=m_{\delta,1}\)。二者均在原点连续，故
\[
 \int_{\mathbb R^d}K_\delta(y)\,dy=m_{\delta,1}(0)=1.
\]
对任何 \(\rho>0\)，缩放后的绝对尾质量满足
\[
 \int_{|x|>\rho}|K_{\delta,R}(x)|\,dx
       =\int_{|y|>R\rho}|K_\delta(y)|\,dy\longrightarrow0.
\]
这些是可以变号的逼近恒等条件。

更直接地，在 \(L^p\) 中有
\[
 B_R^\delta f-f
       =\int K_\delta(y)\,[f(\,\cdot-y/R)-f]\,dy.
\]
有限指数的平移连续性使每个固定 \(y\) 对应的差范数趋零，
且它至多为 \(2\|f\|_p\)。用 \(|K_\delta|\) 作可积控制，
Minkowski 与支配收敛即得范数极限。对有界一致连续输入，
同样使用平移的一致范数连续性，得到一致极限。
\end{proof}

在 \(L^\infty\) 上只陈述有界一致连续输入的收敛，不能把它扩张为
所有有界函数的一致逼近。核卷积可以平滑一个跳跃，却不能在一致范数
中把跳跃的两侧同时逼近到任意精度。后面的习题将把这个边界层算出来。
总积分为一给出平移差公式，绝对值可积性控制误差，因此核允许变号。

\ProseParagraphBreak
可积阈值以上的径向衰减还控制全部半径的极大族。
设 \(\varepsilon=\delta-(d-1)/2>0\)，则
\[
 |K_\delta(x)|\leq C(1+|x|)^{-d-\varepsilon}.
\]
因而本问题满足\cref{b2:prop:17-approximate-identity}的具体核假设。
为看清严格正余量的用途，可以直接将卷积绝对值拆成小球与二进环带：
对 \(R>0\)，小球 \(|y|\leq R^{-1}\) 的贡献至多为 \(CMf(x)\)；
在
\(2^{j-1}R^{-1}<|y|\leq2^jR^{-1}\)、\(j\geq1\)，核大小与球体积相乘后
给出 \(C2^{-j\varepsilon}\)。于是
\[
 |B_R^\delta f(x)|
 \leq C\left(1+\sum_{j\geq1}2^{-j\varepsilon}\right)Mf(x).
\]

\Needspace{6\baselineskip}
\begin{corollary}[极大控制与 Lebesgue 点求和]\label{b2:cor:24-maximal}
设 \(d\geq1\)、\(\delta>(d-1)/2\)、\(1\leq p<\infty\)、
\(f\in L^p(\mathbb R^d)\)。采用核积分的代表，令
\[
 B_*^\delta f(x)\coloneqq\sup_{R>0}|B_R^\delta f(x)|.
\]
则 \(B_*^\delta f\leq C_{d,\delta}Mf\) 几乎处处。
它在 \(p>1\) 时强型有界，在 \(p=1\) 时为弱 \((1,1)\) 型。
在 \(f\) 的每个 Lebesgue 点 \(x\)，有
\[
 B_R^\delta f(x)\longrightarrow f(x)\quad(R\to\infty).
\]
\end{corollary}
\begin{proof}
前面的环带估计在 \(Mf(x)<\infty\) 处同时保证绝对收敛和极大上界。
对固定这样的 \(x\)，半径在任意紧区间内变化时，核函数连续，
并有与区间有关的共同可积控制；可用相同的环带和来作支配收敛。
所以 \(R\mapsto B_R^\delta f(x)\) 连续，上确界可取遍正有理半径，
从而是可测函数。极大函数的强、弱型界给出相应算子界。

若 \(x\) 为 Lebesgue 点，则 \(Mf(x)<\infty\)：小球平均由
Lebesgue 点性质控制，大球平均由 \(L^p\) 范数和 Hölder 控制，
其余半径留在紧区间。将卷积减去 \(f(x)\)，在每个环带上改用
\(|f(x-y)-f(x)|\) 的球平均。固定 \(j\) 时，球半径 \(2^j/R\to0\)，
该平均趋零；所有这些平均又被 \(Mf(x)+|f(x)|\) 控制。
由于 \(\sum_j2^{-j\varepsilon}<\infty\)，先控制级数尾，再令 \(R\to\infty\)，
就得到点态极限。
\end{proof}

如果把 \(\delta\) 降到核可积阈值，环带权不再可和，上述证明停止。
某个固定半径算子仍可能在中间指数上有界，但那并不自动控制极大族。
范数收敛、某个子列的几乎处处收敛与所有半径的几乎处处收敛，
需要各自的依据。

\ProseParagraphBreak
最后把两个端点的常数算清。对任意 \(K\in L^1(\mathbb R^d)\)，
卷积算子在 \(L^1\) 与 \(L^\infty\) 上的准确范数都是 \(\|K\|_1\)。
Young 已给出上界。在 \(L^1\) 上，取非负积分一的光滑近似恒等族
\(\rho_n\)，则
\[
 \|\rho_n\|_1=1,\quad K*\rho_n\to K\text{ 于 }L^1,
\]
这给出反向范数不等式。在 \(L^\infty\) 上，取
\[
 f(-y)=
 \begin{cases}
 \overline{K(y)}/|K(y)|,&K(y)\neq0,\\
 0,&K(y)=0.
 \end{cases}
\]
于是 \(\|f\|_\infty\leq1\) 且 \((K*f)(0)=\|K\|_1\)。
因为 \(K\) 的 \(L^1\) 平移连续，
\[
 |(K*f)(x+h)-(K*f)(x)|
       \leq\|K(\,\cdot+h)-K\|_1\cdot\|f\|_\infty,
\]
输出是连续函数，零点的值也给出本性上确界的下界。

\ProseParagraphBreak
应用于本章的可积核，得到
\[
 \|B_R^\delta\|_{1\to1}
       =\|B_R^\delta\|_{\infty\to\infty}
       =\|K_\delta\|_1>1.
\]
最后一个严格不等式来自振荡渐近：在无穷多径向区间内，
余弦主项为负且支配余项，所以核在一个正测度集合上为负。
核实值且总积分为一，因此
\(\int|K_\delta|=1+2\int(K_\delta)_- >1\)。
热核和 Poisson 核的非负性会给端点压缩；Bochner–Riesz 的总质量
归一化不具有同样的结论。

\ProseParagraphBreak
一维的普通指数范围至此可以完整辨认。零阶区间投影在
\(1<p<\infty\) 上由 Hilbert 变换有界，正阶再由
\cref{b2:prop:24-order-monotonicity}取得同一范围；
\(p=1,\infty\) 则要求正阶，核可积性给出充分性，下一节的尾部检验
给出必要性。高维中，仅靠这些事实仍留下大量中间指数，
这就是 Bochner–Riesz 问题的内容。

% !TeX root = ../../Book2.tex
\section{Bochner–Riesz 问题}
\label{b2:sec:2403}

可积核条件给出同时适用于所有 \(p\) 的范围，却未必给出固定 \(p\) 所需的最小平滑量。中心指数 \(p=2\) 只要求符号有界；离开中心后，核的尾部和不同方向的集中分别产生障碍。本节先给出可以直接检验的必要条件，再证明一个解析插值充分范围，最后在已知乘子有界性的前提下量化逼近速度。必要条件、充分条件和逼近阶数承担不同职责，不能用其中一个代替另外两个。

\subsection{\texorpdfstring{\(L^2\)}{L2} 理论}
\label{b2:subsec:240301}

设 \(d\geq1\)，固定 \(\delta\geq0\)。对每个 \(R>0\)，\cref{b2:def:24-bochner-riesz}中的符号
\(m_{\delta,R}(\xi)=(1-|\xi|^2/R^2)_+^\delta\)
在球内介于零和一之间，在原点邻域任意接近一；零阶时它是球的示性函数。因此符号的本性上确界恰为 \(1\)，\cref{b2:prop:22-l2}给出
\[
 \|B_R^\delta\|_{2\to2}=1,\quad
 \lim_{R\to\infty}\|B_R^\delta f-f\|_2=0
 \quad(f\in L^2).
\]
收敛的具体理由是 \(m_{\delta,R}(\xi)\to1\)，而
\(|m_{\delta,R}-1|^2\cdot|\widehat f|^2\leq|\widehat f|^2\in L^1\)；
先在 Fourier 一侧用支配收敛，再用 Plancherel 返回空间侧。
但 \(L^2\) 范数界不控制频率边界上的点态取值，也不自动控制 \(L^p\)、\(p\ne2\)。与限制问题的联系需要额外结构。

\ProseParagraphBreak
对于一维 \(\delta=0\)，球就是区间，Hilbert 变换给出所有 \(1<p<\infty\) 上的界。高维球不再是有限个半轴投影的乘积。矩形投影可以逐坐标应用 Hilbert 变换并用 Fubini 得到有界性；球边界的法向连续变化，出现了无限方向的组合。

\ProseParagraphBreak
平方理论还允许量化平均逼近的速度，而且显示有无边界消失阶会改变速度结构。
\Needspace{6\baselineskip}
\begin{proposition}\label{b2:prop:24-l2-rate}
设整数 \(d\geq1\)，\(B_R^\delta\) 是 \(L^2(\mathbb R^d;\mathbb C)\) 上符号为 \((1-|\xi|^2/R^2)_+^\delta\) 的 Fourier 乘子，\(R>0\)；零阶符号取开球的示性函数。设 \(\delta>0\)、\(0<s\leq2\)，且
\(f\in L^2(\mathbb R^d;\mathbb C)\) 满足
\((2\pi|\xi|)^s\cdot\widehat f(\xi)\in L^2\)。
以这一符号定义 \((-\Delta)^{s/2}f\)，则
\[
 \|B_R^\delta f-f\|_2
 \leq C_\delta(2\pi R)^{-s}\cdot\|(-\Delta)^{s/2}f\|_2.
\]
若 \(\delta=0\)，则对任意 \(s>0\) 同一结论成立，常数可以取一。
\end{proposition}
\begin{proof}
对 \(a\geq0\)，初等均值估计给出
\[
 |1-(1-a)_+^\delta|\leq\max\{1,\delta\}\min\{1,a\}.
\]
对 \(0<\delta<1\)，在 \(0\leq a\leq1\) 上有
\((1-a)^\delta\geq1-a\)，所以误差不超过 \(a\)；对 \(\delta\geq1\)，同一区间上的导数上界为 \(\delta\)。当 \(a\geq1\)，误差等于一。令 \(a=|\xi|^2/R^2\)。因为 \(0<s\leq2\)，有
\(\min\{1,a\}\leq a^{s/2}\)，所以
\[
 |m_{\delta,R}(\xi)-1|
 \leq C_\delta R^{-s}|\xi|^s.
\]
将两边乘以 \(|\widehat f|\)，取平方积分并用 Plancherel，即得第一项；因 \((-\Delta)^{s/2}\) 的符号是 \((2\pi|\xi|)^s\)，常数中出现 \((2\pi R)^{-s}\)。零阶硬截断的误差只在 \(|\xi|>R\)，在那里
\(\mathbf1_{\{|\xi|>R\}}\leq(|\xi|/R)^s\)
对任意 \(s>0\) 成立，故得到第二项。
\end{proof}
对于 \(\delta>0\)，在一个固定紧频带内有
\(R^2[1-(1-|\xi|^2/R^2)^\delta]\to\delta|\xi|^2\)。
若非零 \(\widehat f\) 支于该频带，支配收敛于是给出
\[
 R^2(B_R^\delta f-f)\longrightarrow
 -\frac{\delta}{4\pi^2}(-\Delta)f
 \quad\text{于 }L^2.
\]
因此通常不能仅靠输入具有更多光滑性把第一项误差提高到 \(R^{-s}\)、\(s>2\)。要提高阶数，应组合不同平均以消去 Taylor 首项，或者采用在低频附近恰好等于一的另一种截断。频率边界的平滑与低频的逼近阶数是两个不同因素。

\ProseParagraphBreak

这个例子也说明“越平滑，逼近必然越快”不是无条件原则。硬截断在固定紧频带上最终误差为零；正阶平均虽然在球边界更光滑，却在球内部偏离一，因而留下低阶 Taylor 误差。硬截断的困难主要是其他 \(L^p\) 的一致有界性和逐点求和；平方逼近本身则十分直接。

\subsection{临界指标}
\label{b2:subsec:240302}

\Needspace{6\baselineskip}
\begin{proposition}[尾部必要条件]\label{b2:prop:24-necessary}
设整数 \(d\geq1\)、\(\delta\geq0\)、\(1\leq p\leq\infty\)，\(B_1^\delta\) 为符号 \((1-|\xi|^2)_+^\delta\) 的 \(L^2\) Fourier 乘子，零阶符号取单位开球的示性函数。若它具有与 Schwartz 输入上的 Fourier 定义一致的有界延拓 \(L^p(\mathbb R^d;\mathbb C)\to L^p(\mathbb R^d;\mathbb C)\)，则在 \(p\ne2\) 时必须满足
\[
 \delta>d\left|\frac1p-\frac12\right|-\frac12.
\]
当右侧为负数时，该条件不提供比 \(\delta\ge0\) 更强的限制。
\end{proposition}
\begin{proof}
选 \(\widehat f\in C_c^\infty\)，在闭单位球附近等于 \(1\)。则 \(f\) 为 Schwartz 函数，且
\[
 B_1^\delta f=K_\delta
 =\mathcal F^{-1}\bigl[(1-|\xi|^2)_+^\delta\bigr].
\]
所以对 \(p<\infty\)，算子有界性迫使 \(K_\delta\in L^p\)。
\cref{b2:prop:24-kernel-tail}给出的精确核可积性条件为
\[
 p\bigl(\delta+(d+1)/2\bigr)>d,
 \quad\text{即}\quad
 \delta>\frac dp-\frac{d+1}{2}.
\]
恢复严格不等式的来源：核主项的余弦在每个足够大的单位径向区间内都有一段固定比例集合保持模下界，余项比主项低一阶。故径向绝对积分至少为下列级数的某个正常数倍：
\[
 \sum_{n\geq n_0}n^{d-1-p(\delta+(d+1)/2)}.
\]
阈值等号时正是发散的调和级数。

符号为实值，在 Schwartz 类上有
\(\langle B_1^\delta f,g\rangle=\langle f,B_1^\delta g\rangle\)。
若 \(1<p<\infty\)，由此和 \(L^p\) 对偶性，\(p\) 上的界也给出 \(p'\) 上相同的界。分别应用刚得的尾部条件，再取较强者，得到
\[
 \delta>\max\left\{\frac dp,\frac d{p'}\right\}
 -\frac{d+1}{2}
 =d\left|\frac1p-\frac12\right|-\frac12.
\]

若 \(p=1\)，直接得到 \(\delta>(d-1)/2\)。若 \(p=\infty\)，对 Schwartz 函数 \(f,g\) 使用实符号配对和假设的 \(L^\infty\) 界，有
\[
 |\langle B_1^\delta f,g\rangle|
 =|\langle f,B_1^\delta g\rangle|
 \leq C\|f\|_1\cdot\|g\|_\infty.
\]
对紧支撑光滑 \(g\) 的单位球取上确界，便得到
\(\|B_1^\delta f\|_1\leq C\|f\|_1\)，于是也可应用 \(p=1\) 的必要条件。给定延拓在 Schwartz 类上的配对关系与 \(L^1\) 范数的对偶公式，便把无穷指数的问题转到可积端点。对于 \(p=2\)，尾部检验给出 \(\delta>-1/2\)，所有 \(\delta\geq0\) 都满足。
\end{proof}
严格性来自径向调和级数；没有对尾部作精确下界，就不能宣称阈值等号也失败。这个检验并未排除临界函数
\(\max\{d|1/p-1/2|-1/2,0\}\) 中为零的内部指数上的硬截断；下一小节的深层几何定理才排除高维硬球。

\subsection{Bochner–Riesz 猜想}
\label{b2:subsec:240303}

经典 Bochner–Riesz 目标是：对 \(d\ge2\)、\(1<p<\infty\)，在
\[
 \delta>\max\left\{d\left|\frac1p-\frac12\right|-\frac12,\ 0\right\}
\]
时取得 \(L^p\) 有界性。这里严格的正平滑量将硬球情形排除在目标外。维数、端点形式、径向输入限制以及局部版本各是不同问题。

\ProseParagraphBreak

下面采用\term[fefferman-qiuchengzi-dingli]{Fefferman 球乘子定理}[Fefferman ball multiplier theorem]的下述版本。

\Needspace{6\baselineskip}
\begin{theorem}[深层结构定理：Fefferman 球乘子定理]\label{b2:thm:24-fefferman}
设整数 \(d\geq2\)，\(1<p<\infty\)。在 \(L^2(\mathbb R^d;\mathbb C)\) 上以符号
\(\mathbf1_{\{\,|\xi|<1\,\}}\) 定义单位球 Fourier 乘子。它具有与 Schwartz 类上一致的有界延拓
\(L^p(\mathbb R^d;\mathbb C)\to L^p(\mathbb R^d;\mathbb C)\)，当且仅当 \(p=2\)。
\end{theorem}
\begin{proofsketch}
假设球乘子在某个非二指数有界。对球作放大、频移并取极限，配合随机符号，可推出各方向半空间投影的一致向量值界。Besicovitch 型细矩形构造使输入的平方函数质量分散，而输出在一个小测度集合中累积；对 \(p>2\) 这与向量值界矛盾。对偶性处理 \(p<2\)，切片传递处理高维。
几何矩形构造及传递论证的详细证明见原文
\cite[Theorem 1; Section 2, pp.~330--336]{Fefferman1971Ball}。
\end{proofsketch}
这个定理并不说任何正的 \(\delta\) 都足够。球乘子的失败与正平滑后的范围之间，仍隔着一个依赖维数与 \(p\) 的估计问题。

\subsection{解析插值与 \texorpdfstring{\(L^p\)}{Lᵖ} 逼近估计}
\label{b2:subsec:240304}

可积核给出了 \(L^1\) 和 \(L^\infty\) 两个强端点，平方理论给出了中心指数。它们对应不同的平滑阶数；要连接它们，必须同时改变算子中的参数，而不能只对一个未知有界的固定算子插值。先证明复参数竖线上的核估计，再把参数与输入指数联系起来。

\Needspace{6\baselineskip}
\begin{lemma}[固定实部的复阶 Bochner--Riesz 核]\label{b2:lem:24-complex-kernel}
设整数 \(d\geq1\)，固定实数 \(\sigma>-1\)。对 \(z=\sigma+\mathrm it\)、\(t\in\mathbb R\)，令
\[
 m_z(\xi)=
 \begin{cases}
  \exp\bigl(z\log(1-|\xi|^2)\bigr),&|\xi|<1,\\
  0,&|\xi|\geq1,
 \end{cases}
 \quad K_z=\mathcal F^{-1}m_z,
\]
其中球内使用实对数。符号 \(m_z\) 属于 \(L^1\)，所以 \(K_z\) 是有界连续函数。存在只依赖 \(d,\sigma\) 的常数 \(C\)，使
\[
 |K_{\sigma+\mathrm it}(x)|
 \leq C e^{C(1+t^2)}
 (1+|x|)^{-\sigma-(d+1)/2},
 \quad x\in\mathbb R^d,\ t\in\mathbb R.
\]
特别地，若 \(\sigma>(d-1)/2\)，则
\(\|K_{\sigma+\mathrm it}\|_1\leq C e^{C(1+t^2)}\)。
\end{lemma}
\begin{proof}
记 \(\nu=d/2+z=u+\mathrm it\)，其中 \(u=d/2+\sigma>-1/2\) 固定。\subsecref{b2:subsec:230402}已经以整函数 \(P=1/\Gamma\) 构造复阶 Bessel 级数，并证明归一化球内分布的核恒等式。这里 \(\operatorname{Re}z>-1\)，是普通可积符号，乘回 Gamma 因子即得
\[
 K_z(x)=\Gamma(z+1)\pi^{-z}
 |x|^{-\nu}J_\nu(2\pi|x|),\quad x\ne0.
\]
这是未归一化平均的核，因子 \(\Gamma(z+1)\) 必须保留。复 Gamma 的积分定义给出
\[
 |\Gamma(\sigma+1+\mathrm it)|
 \leq\int_0^\infty e^{-v}v^\sigma\,\mathrm dv
 =\Gamma(\sigma+1).
\]
因此分子 Gamma 因子在这条竖线上甚至一致有界。

小半径与初值由级数控制。对 \(0<r\leq1\)，
\[
 r^{-\nu}J_\nu(2\pi r)
 =\pi^\nu\sum_{k=0}^\infty
 \frac{(-\pi^2r^2)^k}{k!\Gamma(k+\nu+1)}.
\]
因为 \(u+1>0\)，Gamma 递推给出
\[
 |\Gamma(k+\nu+1)|^{-1}
 \leq
 \frac{|P(u+1+\mathrm it)|}
 {\prod_{\ell=0}^{k-1}(u+1+\ell)}.
\]
分母是增长的正实乘积；配合 \(k!\)，级数及其任意固定阶 \(r\) 导数均可在 \(0\leq r\leq1\) 一致求和。前章从 Gamma 无限乘积得到
\[
 |P(u+1+\mathrm it)|\leq C_u e^{C_u(1+t^2)}.
\]
于是 \(|r^{-\nu}J_\nu(2\pi r)|\leq C e^{C(1+t^2)}\)，并且
\[
 |J_\nu(2\pi)|+
 \left|\frac{\mathrm d}{\mathrm dr}J_\nu(2\pi r)
 \bigg|_{r=1}\right|
 \leq C e^{C(1+t^2)}.
\]
对 \(r^\nu\) 求导带来的 \(1+|\nu|\) 因子也可并入此界。前面的 Gamma 分子界因而给出 \(0\leq|x|\leq1\) 的目标估计，原点按连续极限解释。

对 \(r\geq1\)，令
\[
 w(r)=r^{1/2}J_\nu(2\pi r),\quad
 E(r)=|w'(r)|^2+4\pi^2|w(r)|^2.
\]
复阶 Bessel 方程仍给出
\[
 w''+4\pi^2w=(\nu^2-1/4)r^{-2}w.
\]
能量 \(E\) 是实的非负量，不需要系数 \(\nu^2\) 为实数。微分并消去自由振荡项后，
\[
 |E'(r)|
 \leq C|\nu^2-1/4|r^{-2}\cdot E(r)
 \leq C_u(1+t^2)r^{-2}\cdot E(r).
\]
积分 \(\int_1^\infty r^{-2}\,\mathrm dr=1\)，结合已证的 \(r=1\) 初值界，得到
\[
 E(r)\leq E(1)e^{C_u(1+t^2)}
 \leq C e^{C(1+t^2)},\quad r\geq1.
\]
因此
\[
 |J_\nu(2\pi r)|
 \leq C e^{C(1+t^2)}r^{-1/2}.
\]
核公式中另外的幂因子模为 \(\pi^{-\sigma}r^{-u}\)，所以大半径衰减为
\(r^{-u-1/2}=r^{-\sigma-(d+1)/2}\)。
与小半径估计合并即得结论。最后，当
\(\sigma+(d+1)/2>d\) 时，径向积分可积，便得到统一的 \(L^1\) 核界。
\end{proof}

与球面限制中的竖线不同，此处 \(\nu\) 的实部是 \(d/2+\sigma\)。能量法仍然适用，因为对固定实部有
\(|\nu^2-1/4|\leq C(1+t^2)\)；不能直接把只对固定实阶陈述的 Bessel 渐近当作这一统一估计。另一个区别在插值的目标：本节要得到 \(L^p\to L^p\)，而\cref{b2:lem:23-compact-analytic-interpolation}的陈述给出的是 \(L^p\to L^{p'}\)。下面保留其共同定义域的做法，重新选择测试函数指数。

\Needspace{6\baselineskip}
\begin{theorem}[可直接证明的插值范围]\label{b2:thm:24-interpolation-range}
设整数 \(d\geq1\)、\(1<p<\infty\)，实数 \(\delta\) 满足
\[
 \delta>(d-1)\left|\frac1p-\frac12\right|,
\]
并在 Schwartz 类上以
\[
 B_R^\delta f=\mathcal F^{-1}
 \bigl[(1-|\xi|^2/R^2)_+^\delta\cdot\widehat f(\xi)\bigr],
 \quad R>0,
\]
定义 Bochner--Riesz 平均。则这些算子具有与 \(L^2\) 定义相容的唯一的 \(L^p(\mathbb R^d;\mathbb C)\) 有界延拓，且
\[
 \sup_{R>0}\|B_R^\delta\|_{p\to p}\leq C_{d,p,\delta},
 \quad
 \lim_{R\to\infty}\|B_R^\delta f-f\|_p=0
 \quad(f\in L^p).
\]
\end{theorem}
\begin{proof}
当 \(p=2\) 时，结论已由本节开头的 Plancherel 论证得到，甚至允许 \(\delta=0\)。先设 \(1<p<2\)，令
\[
 \theta=\frac2p-1\in(0,1),\quad
 \sigma=\frac\delta\theta.
\]
题设恰好使 \(\sigma>(d-1)/2\)，而目标参数满足 \(\delta=\theta\sigma\)。在闭条带 \(0\leq\operatorname{Re}w\leq1\) 上，考察
\[
 S_wf=\mathcal F^{-1}(m_{\sigma w}\cdot\widehat f),
 \quad
 m_{\sigma w}(\xi)=(1-|\xi|^2)_+^{\sigma w}.
\]
球内使用实对数，球外为零，球面零测集的取值无关。在整条带上
\(|m_{\sigma w}|\leq1\)，且符号支撑于闭单位球。
对内点参数的任意固定阶导数，只多出
\(\sigma^k\log^k(1-|\xi|^2)\)；
在参数紧集上，其绝对积分由球面法向变量中可积的
\(u^\varepsilon|\log u|^k\) 控制。因此符号与光滑频率测试函数的配对全纯，并连续到两条边界。

先明确一套可随参数变化的共同测试空间。令 \(f,g\) 是有界紧支撑简单函数。它们的 Fourier 变换有界，所以
\[
 B_w(f,g)=\int_{|\xi|<1}
 m_{\sigma w}(\xi)\cdot\widehat f(\xi)
 \cdot\widehat g(-\xi)\,\mathrm d\xi
\]
是绝对收敛的双线性配对。它也等于
\(\int(S_wf)(x)\cdot g(x)\,\mathrm dx\)：符号可积使核有界，两个输入紧支撑，故两种写法由 Fubini 相等。这个共同定义域包括稍后使用的所有变化测试函数。

取 \(\|f\|_p=\|g\|_{p'}=1\)，把它们的非零值分别写成
\(f=\sum_j a_j\mathbf1_{E_j}\)、
\(g=\sum_k b_k\mathbf1_{F_k}\)，同一族内的集合两两不交且有界。定义
\[
 \begin{aligned}
 f_w&=\sum_j\frac{a_j}{|a_j|}
 |a_j|^{p(1+w)/2}\cdot\mathbf1_{E_j},\\
 g_w&=\sum_k\frac{b_k}{|b_k|}
 |b_k|^{p'(1-w)/2}\cdot\mathbf1_{F_k}.
 \end{aligned}
\]
在 \(w=\theta\) 时，两个复幂指数都等于一，所以
\(f_\theta=f\)、\(g_\theta=g\)。在竖边上则有
\[
 \|f_{\mathrm it}\|_2=\|g_{\mathrm it}\|_2=1,\quad
 \|f_{1+\mathrm it}\|_1=1,\quad
 \|g_{1+\mathrm it}\|_\infty=1.
\]
测试函数 \(g_w\) 的指数与 \(f_w\) 方向相反，正是为了在右边界配对 \(L^1\) 输出与 \(L^\infty\) 测试。

函数 \(F(w)=B_w(f_w,g_w)\) 在条带内全纯并连续到边界，因为 \(f_w,g_w\) 都是有限解析组合。并且各系数的模仅依赖 \(\operatorname{Re}w\)，于是
\[
 |F(w)|\leq |B(0,1)|\cdot
 \|f_w\|_1\cdot\|g_w\|_1\leq C_{f,g}
\]
在整个闭条带成立。有限解析组合及这项统一上界共同保证了条带插值所需的解析性与增长控制。

在 \(w=\mathrm it\) 上，符号模不超过一，Plancherel 与 Cauchy--Schwarz 给出
\(|F(\mathrm it)|\leq1\)。
在 \(w=1+\mathrm it\) 上，\cref{b2:lem:24-complex-kernel}及 Young 不等式给出
\[
 \|S_{1+\mathrm it}\|_{1\to1}
 \leq C e^{C(1+\sigma^2t^2)},
 \quad
 |F(1+\mathrm it)|\leq C e^{C(1+\sigma^2t^2)}.
\]
取固定 \(\kappa>C\sigma^2\)，对
\(e^{\kappa(w-\theta)^2}F(w)\)
应用\cref{b2:thm:17-three-lines}。这个因子在竖直方向压住右边界增长，在插值点等于一，边界的实部只带来固定常数。由此得到
\[
 |B_\theta(f,g)|\leq C_{d,p,\delta}.
\]
齐次性与 \(L^p\) 对偶性说明
\(\|S_\theta f\|_p\leq C_{d,p,\delta}\|f\|_p\)
在紧支撑简单函数上成立；稠密性给出唯一的 \(L^p\) 延拓。
它与原 Fourier 定义相容：对 Schwartz 输入作有界紧支撑简单函数逼近，可同时在 \(L^1,L^2,L^p\) 中收敛，频率配对与 \(L^2\) 乘子极限均保持一致。因此 \(S_\theta=B_1^\delta\)。

若 \(2<p<\infty\)，则 \(1<p'<2\)，并且
\(|1/p'-1/2|=|1/p-1/2|\)。
实值符号 \(m_\delta\) 使 \(B_1^\delta\) 在测试类上自伴随，已证的 \(L^{p'}\) 界通过对偶性给出 \(L^p\) 界。最后，\cref{b2:prop:24-dilation}的同指数共轭公式使各半径范数相同，\cref{b2:prop:24-strong-approximation}用统一算子界与稠密测试类给出 \(R\to\infty\) 时的强收敛。
\end{proof}
这是一个明确已证范围，一般窄于经典目标。它的证明只结合 \(L^2\) 与可积核端点，没有利用球面限制的额外改进。比较两条临界直线，能看清深层理论应在哪个区间增加信息。

\ProseParagraphBreak
已经证明的一致乘子界也能与输入正则性配合，给出其他 \(L^p\) 中的逼近速度。
\Needspace{6\baselineskip}
\begin{proposition}\label{b2:prop:24-lp-rate}
设整数 \(d\geq1\)，\(1<p<\infty\)，\(\delta>0\)，\(0<s\leq2\)。设符号 \((1-|\xi|^2/R^2)_+^\delta\) 的 Fourier 乘子具有与 \(L^2\) 定义相容的延拓 \(B_R^\delta:L^p\to L^p\)，并满足
\[
 M\coloneqq\sup_{R>0}\|B_R^\delta\|_{p\to p}<\infty.
\]
令 \(H^{s,p}(\mathbb R^d)\) 按\cref{b2:def:22-bessel-space}定义，即
\(\|f\|_{H^{s,p}}=\|J^sf\|_p\)，其中 \(J^s\) 的符号为
\((1+4\pi^2|\xi|^2)^{s/2}\)。则对每个 \(f\in H^{s,p}(\mathbb R^d)\) 和 \(R\geq1\)，
\[
 \|B_R^\delta f-f\|_p
 \leq C_{d,p,s,\delta,M}R^{-s}\cdot\|f\|_{H^{s,p}}.
\]
\end{proposition}
\begin{proof}
先对 Schwartz 输入证明；该类在 \(H^{s,p}\) 中稠密，且 \(J^{-s}\) 在 \(L^p\) 上有界，最终可把估计传到一般输入。取固定径向
\(\chi\in C_c^\infty(\mathbb R^d)\)，在 \(|\xi|\leq1/4\) 上为一，支撑于 \(|\xi|<1/2\)。令
\(P_Rf=\mathcal F^{-1}[\chi(\xi/R)\cdot\widehat f]\)，并记 \(g=J^sf\)。分解
\[
 (I-B_R^\delta)f
 =(I-B_R^\delta)P_Rf+
 (I-B_R^\delta)(I-P_R)f.
\]
后一个输入只含 \(|\xi|\geq R/4\) 的高频。把其 \(R^{-s}\) 因子提出后，作用于 \(g\) 的乘子为
\[
 a_R(\xi)=R^s(1-\chi(\xi/R))
 (1+4\pi^2|\xi|^2)^{-s/2}.
\]
在其支撑上 \(R\leq4|\xi|\)。每次微分 Bessel 因子带来相应的 \((1+|\xi|)^{-1}\)，微分截断的项只出现在 \(|\xi|\asymp R\) 上。因此对 \(|\alpha|\leq d+2\)，
\[
 |\xi|^{|\alpha|}\cdot|\partial^\alpha a_R(\xi)|
 \leq C_\alpha R^s(1+|\xi|)^{-s}\leq C_\alpha
\]
一致成立。\cref{b2:thm:22-mihlin}给出
\(\|(I-P_R)f\|_p\leq C R^{-s}\|g\|_p\)，再应用
\(\|I-B_R^\delta\|_{p\to p}\leq1+M\)
便控制高频误差。

低频误差的归一化乘子为
\[
 \ell_R(\xi)=R^s[1-(1-|\xi|^2/R^2)^\delta]\cdot
 \chi(\xi/R)(1+4\pi^2|\xi|^2)^{-s/2}.
\]
在截断支撑上 \(v=|\xi|^2/R^2\leq1/4\)，远离符号的球边界。由微积分基本定理，
\[
 1-(1-v)^\delta=v\,b_\delta(v),\quad
 b_\delta(v)=\delta\int_0^1(1-tv)^{\delta-1}\,\mathrm dt.
\]
所以 \(b_\delta\) 及每个所需阶数的导数在 \(0\leq v\leq1/4\) 有界。
把低频乘子写成
\[
 \ell_R(\xi)=R^{s-2}|\xi|^2\cdot
 b_\delta(|\xi|^2/R^2)\cdot\chi(\xi/R)
 (1+4\pi^2|\xi|^2)^{-s/2}.
\]
链式与乘积求导可分成两个区域记录：对所有 \(|\alpha|\leq d+2\)，
\[
 |\xi|^{|\alpha|}\cdot|\partial^\alpha\ell_R(\xi)|
 \leq C_\alpha
 \begin{cases}
 R^{s-2}|\xi|^2,&0<|\xi|\leq1,\\
 (|\xi|/R)^{2-s},&1\leq|\xi|\leq R/2.
 \end{cases}
\]
第一行利用了 \(|\xi|^2\) 的零点和其余因子在单位球内的光滑有界性；高阶导数即使消去这个零点，外面的 \(|\xi|^{|\alpha|}\) 也补足所需幂次。第二行中，微分 \(b_\delta\) 或截断产生的尺度因子由 \(R^{-1}\lesssim|\xi|^{-1}\) 控制，微分 Bessel 因子带来 \(|\xi|^{-1}\)。这解释了每次求导后仍保持同一零阶界的原因。

因为 \(R\geq1\)、\(s\leq2\)，两行都一致有界；支撑外乘子为零。再用 Mihlin 定理，得到低频误差不超过 \(C R^{-s}\|g\|_p\)。合并两部分后，以 \(H^{s,p}\) 中的 Schwartz 逼近传到一般 \(f\)，即得结论。
\end{proof}
在一致乘子界成立的指数范围内，这个论证把两种误差分开：高频尾部由输入的光滑性控制，低频误差由平均符号在原点的 Taylor 阶数控制。若进一步构造消掉低阶 Taylor 系数的组合平均，后一限制可以改变，但需要重新核对组合的统一算子界。

\ProseParagraphBreak

例如，对固定 \(d,p\) 选择处于已证插值范围的 \(\delta\)，由\cref{b2:prop:22-integer-sobolev}的整数阶范数等价，一个 \(W^{1,p}\) 输入有 \(R^{-1}\) 的范数逼近速度，一个 \(W^{2,p}\) 输入有 \(R^{-2}\) 的速度。这里只是范数误差，不能把它直接变成全参数逐点误差率；后者需要位置上的极大控制或更强的正则性与嵌入。

\ProseParagraphBreak

必要直线与插值直线的距离可以用一个数值模型说明。在 \(d=3\)、\(p=4\) 中，尾部条件要求 \(\delta>1/4\)，本章直接证明的插值范围要求 \(\delta>1/2\)。例如 \(\delta=3/4\) 在已证范围内，\(\delta=1/8\) 被必要条件排除，而 \(\delta=3/8\) 不能仅靠这两项结果判断。说它“没有被反例排除”与说它“已证明有界”有明确区别。

\ProseParagraphBreak

指数对偶把 \(p=4\) 与 \(p=4/3\) 的范围对应起来，因为球符号为实且算子在测试类上自伴随。缩放又使每个已证同指数界自动对所有半径具有同一常数。因此进一步扩大范围的关键，是获得更精细的平滑量与输入指数关系。若一篇文献只给固定空间球内的局部输出或径向输入，其命题还不能直接填入这个全空间、一般输入的空缺。

\ProseParagraphBreak

在 \(p\) 接近二时，尾部必要条件的右侧可能为负，此时它只要求原来已经约定的 \(\delta\geq0\)，看起来没有门槛。然而高维 \(\delta=0\) 已由球乘子定理排除所有非二指数。这个额外的几何障碍来自多方向叠加，并不是单个核尾部的一条更仔细积分能看出来的。正阶问题是否能在接近二的范围内降到任意小的平滑量，则要依赖更深的曲率和限制估计。

\ProseParagraphBreak

对角线上的 \(p=2\) 则没有这些困难：符号有界已经给出范数一，且所有 \(\delta\geq0\) 都随固定阶数的半径增加强收敛到恒等。这个特殊点是各种插值证明的可靠起点，但它不控制其他指数上的几何组合。把一个证明完全归结为“平方理论加插值”时，应能够说清另一个端点来自哪里，以及为什么它有足够强的范围和统一常数。
\ProseParagraphBreak
消去低频误差的组合可以直接构造，而无需增加单个平均的边界消失阶。半径放大到 \(\sqrt2R\) 时，低频二次项的系数减半；用权 \(2,-1\) 相减，就能保留常数项并消去该误差。

\Needspace{6\baselineskip}
\begin{proposition}[两半径组合的逼近阶数]\label{b2:prop:24-two-radius-rate}
设整数 \(d\geq1\)，\(1<p<\infty\)，\(\delta>0\)，且符号
\((1-|\xi|^2/R^2)_+^\delta\) 的相容 \(L^p\) Fourier 乘子满足
\(\sup_{R>0}\|B_R^\delta\|_{p\to p}\leq M<\infty\)。
定义
\[
 C_R=2B_{\sqrt2R}^\delta-B_R^\delta,\quad R>0.
\]
则 \(\|C_R\|_{p\to p}\leq3M\)，并且对 \(0<s\leq4\)、
\(R\geq1\) 和 \(f\in H^{s,p}(\mathbb R^d)\)，有
\[
 \|C_Rf-f\|_p
 \leq C_{d,p,s,\delta,M}R^{-s}\cdot\|f\|_{H^{s,p}}.
\]
这里 \(H^{s,p}\) 及其范数按\cref{b2:def:22-bessel-space}定义。
若 \(\delta=1\)，上述速率对每个 \(s>0\) 都成立，常数可依赖 \(s\)。
\end{proposition}
\begin{proof}
范数界由两个平均的界直接相加得到。对低频区域 \(|\xi|<R/2\)，令 \(v=|\xi|^2/R^2\)，则组合符号为
\[
 c_\delta(v)=2(1-v/2)^\delta-(1-v)^\delta
 =1-\frac{\delta(\delta-1)}4v^2+O(v^3).
\]
更精确地，\(c_\delta(0)=1\)、\(c_\delta'(0)=0\)，所以 Taylor 积分余项把 \(1-c_\delta(v)\) 写成
\(v^2b(v)\)，其中 \(b\) 在 \([0,1/4]\) 连同所有所需导数有界。

使用\cref{b2:prop:24-lp-rate}证明中的同一低频截断 \(P_R\)，并记 \(g=J^sf\)。高频输入估计
\[
 \|(I-P_R)f\|_p\leq C R^{-s}\cdot\|g\|_p
\]
实际上对任意 \(s>0\) 成立，因为该处仅使用 \(|\xi|\geq R/4\)，不使用 \(s\leq2\)。对它应用 \(I-C_R\)，代价至多为 \(1+3M\)。

低频归一化误差乘子现在是
\[
 R^{s-4}|\xi|^4\cdot b(|\xi|^2/R^2)\cdot
 \chi(\xi/R)(1+4\pi^2|\xi|^2)^{-s/2}.
\]
对 \(|\alpha|\leq d+2\)，其第 \(\alpha\) 阶导数乘以 \(|\xi|^{|\alpha|}\) 后，在 \(0<|\xi|\leq1\) 上由
\(C_\alpha R^{s-4}|\xi|^4\) 控制，在
\(1\leq|\xi|\leq R/2\) 上由
\(C_\alpha(|\xi|/R)^{4-s}\) 控制。理由与前一命题相同：原点的四阶零点控制小频率，归一化链式求导控制大频率。因为 \(0<s\leq4\)、\(R\geq1\)，这些界与 \(R\) 无关。Mihlin 定理给出低频误差的 \(R^{-s}\) 界，再合并高频部并由稠密性推广。

若 \(\delta=1\)，则
\[
 2(1-v/2)-(1-v)=1
\]
在整个低频区域精确成立，低频误差为零；只剩已经对任意 \(s>0\) 成立的高频估计。这证明最后一个断言，但仍使用了一阶平均原先具有统一 \(L^p\) 界的假设。
\end{proof}

组合权有负数，所以它不保留正性，也不能由正权阶数单调公式读出。它改变的是球内部的逼近阶数，而不是每条球边界的曲率障碍。特别是 \(\delta=1\) 时，精确低频抵消改善逼近阶数的作用，仍限于该乘子有界的指数范围。

\ProseParagraphBreak
通常四阶速率也会饱和。取非零 Schwartz 函数 \(f\)，使 \(\widehat f\) 光滑紧支撑，并固定 \(\delta>0\)、\(\delta\ne1\)。当 \(R\) 足够大时，整个频率支撑都在低频区；上面的 Taylor 展开连同频率导数一致成立，所以逆变换在 Schwartz 拓扑中给出
\[
 R^4(C_R-I)f
 \longrightarrow
 -\frac{\delta(\delta-1)}{64\pi^4}(-\Delta)^2f.
\]
系数来自 \(v^2=|\xi|^4/R^4\) 与
\((-\Delta)^2\) 的符号 \(16\pi^4|\xi|^4\)。该极限非零，因此不能只要求输入再光滑一些，就把一般两半径组合的阶数提高到 \(s>4\)。若要继续消去误差，必须增加组合条件或改变平均；稳定的算子界与低频逼近阶数应分别设计。

% !TeX root = ../../Book2.tex
\section{与限制问题的关系}
\label{b2:sec:2404}

球面限制控制薄壳上的 Fourier 质量，Bochner--Riesz 符号则把不同球面层加权求和。两者共享曲率与频率集中，但限制估计首先产生输入、输出指数不同的界。本节把这条联系中的测度与缩放写清，再用空间局部化和多面体逼近说明：从混合指数估计到同指数乘子估计，还缺少哪些控制。

\subsection{球面 Fourier 变换}
\label{b2:subsec:240401}

先取 \(d\geq2\)，单位球面 \(S^{d-1}\) 的面积测度记为 \(\sigma\)，半径 \(r>0\) 的球面面积测度记为 \(\sigma_r\)。它们满足
\[
 \int_{rS^{d-1}}h(\xi)\,\mathrm d\sigma_r(\xi)
 =r^{d-1}\int_{S^{d-1}}h(r\omega)\,\mathrm d\sigma(\omega).
\]
右侧的 \(r^{d-1}\) 来自曲面面积的缩放，与径向积分中的 \(\mathrm dr\) 分别保留。

对 \(f\in\mathcal S(\mathbb R^d)\)，定义单位参数下的球面层算子
\[
 T_rf(x)\coloneqq
 \int_{S^{d-1}}e^{2\pi\mathrm i r x\cdot\omega}
       \cdot\widehat f(r\omega)\,\mathrm d\sigma(\omega).
\]
以真实面积测度积分的限制再延拓算子则是 \(r^{d-1}T_r\)。因此极坐标给出
\begin{equation}\label{b2:eq:24-layer-measure}
 \begin{aligned}
 B_1^\delta f(x)
 &=\int_0^1(1-r^2)^\delta
       \int_{rS^{d-1}}e^{2\pi\mathrm i x\cdot\xi}
          \cdot\widehat f(\xi)\,\mathrm d\sigma_r(\xi)\,\mathrm dr\\
 &=\int_0^1(1-r^2)^\delta r^{d-1}T_rf(x)\,\mathrm dr,
 \quad \delta>-1.
 \end{aligned}
\end{equation}
对 \(\delta\geq0\)，这就是\cref{b2:def:24-bochner-riesz}中的平均；对 \(-1<\delta<0\)，此处先按球内可积符号的逆 Fourier 积分解释。测试函数的 Fourier 变换有界，而
\(\int_0^1(1-r^2)^\delta r^{d-1}\,\mathrm dr<\infty\)，所以换序绝对合法。这个公式中若使用 \(\sigma_r\)，就不应再另乘一次 \(r^{d-1}\)。

\ProseParagraphBreak
内层算子只保留一个频率半径，径向积分才恢复球内部的频率。球面 Fourier 变换的振荡尾部提供 \((d-1)/2\) 阶衰减；径向权在 \(r=1\) 的消失阶又在端点振荡积分中提供 \(\delta+1\) 阶增益，合起来与本章核的 \(\delta+(d+1)/2\) 阶衰减一致。若先对每一层取绝对值，再对半径求和，不同半径之间的抵消便不再可用。下面先看这种直接估计究竟能够证明什么。

\subsection{限制估计}
\label{b2:subsec:240402}

本处需要\cref{b2:thm:23-stein-tomas}的球面端点版本。设
\[
 p_*=\frac{2(d+1)}{d+3},\quad
 q_*=p_*'=\frac{2(d+1)}{d-1}.
\]
对单位球面及面积测度，限制算子 \(Rf=\widehat f|_{S^{d-1}}\) 满足
\[
 \|Rf\|_{L^2(S^{d-1})}\leq A_d\|f\|_{p_*}.
\]
其有界伴随为
\[
 Eg(x)=\int_{S^{d-1}}e^{2\pi\mathrm i x\cdot\omega}
        \cdot g(\omega)\,\mathrm d\sigma(\omega),\quad
 \|Eg\|_{q_*}\leq A_d\|g\|_2.
\]
这里对一般 \(L^{p_*}\) 输入的限制按连续延拓定义，并不是任意 Fourier 代表的逐点取值。球面为有限正测度，故\cref{b2:prop:23-ttstar}也给出 \(ER:L^{p_*}\to L^{q_*}\) 及其范数上界 \(A_d^2\)。

\ProseParagraphBreak
半径改变时，需要同时缩放输入和输出。令 \(f_r(y)=f(y/r)\)，则
\[
 \widehat{f_r}(\omega)=r^d\widehat f(r\omega),\quad
 \|f_r\|_{p_*}=r^{d/p_*}\|f\|_{p_*}.
\]
因此
\[
 \|\widehat f(r\,\cdot)\|_{L^2(S^{d-1})}
 \leq A_dr^{-d/q_*}\|f\|_{p_*},
\]
\[
 T_rf(x)=r^{-d}(ERf_r)(rx),\quad
 \|T_rf\|_{q_*}
 \leq A_d^2r^{-2d/q_*}\|f\|_{p_*}.
\]
最后一个幂也等于 \(d/p_*-d/q_*-d\)。在\eqref{b2:eq:24-layer-measure}中乘以真实面积因子 \(r^{d-1}\) 后，层算子的界便带有
\[
 r^{d-1-2d/q_*}=r^\alpha,\quad
 \alpha\coloneqq\frac{d-1}{d+1}.
\]
这一幂次在零半径处仍然可积，稍后将用于全平均。

\ProseParagraphBreak
把逐半径估计先用于非负权，可以同时得到频率平方界与混合指数界。

\Needspace{6\baselineskip}
\begin{proposition}[非负径向权的限制界]\label{b2:prop:24-radial-weight}
设整数 \(d\geq2\)，
\[
 p_*=\frac{2(d+1)}{d+3},\quad
 q_*=\frac{2(d+1)}{d-1},\quad
 \alpha=\frac{d-1}{d+1}.
\]
设 \(w:(0,\infty)\to[0,\infty)\) 可测且
\[
 I_w\coloneqq\int_0^\infty w(r)r^\alpha\,\mathrm dr<\infty.
\]
对 \(f\in\mathcal S(\mathbb R^d)\)，定义
\[
 A_wf(\xi)=w(|\xi|)^{1/2}\cdot\widehat f(\xi),\quad
 S_wf=\mathcal F^{-1}\!\left[w(|\xi|)\cdot\widehat f(\xi)\right].
\]
则这两个算子分别唯一连续延拓为
\[
 A_w:L^{p_*}(\mathbb R^d)\longrightarrow L^2(\mathbb R^d),\quad
 S_w:L^{p_*}(\mathbb R^d)\longrightarrow L^{q_*}(\mathbb R^d),
\]
并有
\[
 \|A_wf\|_2\leq A_d I_w^{1/2}\|f\|_{p_*},\quad
 \|S_wf\|_{q_*}\leq A_d^2 I_w\|f\|_{p_*}.
\]
常数 \(A_d\) 可取单位球面 Stein--Tomas 限制界的常数，与 \(w,f\) 无关。
\end{proposition}
\begin{proof}
先取 Schwartz 输入。极坐标、非负性和刚才的半径限制估计给出
\[
 \begin{aligned}
 \|A_wf\|_2^2
 &=\int_0^\infty w(r)r^{d-1}
       \|\widehat f(r\,\cdot)\|_{L^2(S^{d-1})}^2\,\mathrm dr\\
 &\leq A_d^2\|f\|_{p_*}^2
       \int_0^\infty w(r)r^{d-1-2d/q_*}\,\mathrm dr
 =A_d^2 I_w\|f\|_{p_*}^2.
 \end{aligned}
\]
这就给出第一项连续延拓。第二项的测试函数定义也有意义：在 \(r\leq1\) 时 \(r^{d-1}\leq r^\alpha\)；在 \(r\geq1\) 时，\(\widehat f\) 的任意阶衰减控制 \(r^{d-1-\alpha}\)。所以 \(w(|\xi|)\cdot\widehat f(\xi)\) 绝对可积。

将 \(S_wf\) 与另一 Schwartz 函数 \(h\) 配对，Fubini 给出
\[
 \langle S_wf,h\rangle
 =\int_{\mathbb R^d}w(|\xi|)\cdot
       \widehat f(\xi)\overline{\widehat h(\xi)}\,\mathrm d\xi
 =\langle A_wf,A_wh\rangle.
\]
于是
\[
 |\langle S_wf,h\rangle|
 \leq A_d^2 I_w\|f\|_{p_*}\cdot\|h\|_{p_*}.
\]
\(L^{p_*}\) 的对偶为 \(L^{q_*}\)，所以这个分布由一个 \(L^{q_*}\) 函数表示，并得到第二项估计。最后由 Schwartz 类在 \(L^{p_*}\) 中稠密，完成唯一连续延拓。
\end{proof}

这里的平方配对就是 \(TT^*\) 方法，但已经在实际使用的径向权上重新给出了论证，未要求这个加权测度具有紧支集。权的非负性保证 \(\langle S_wf,f\rangle\geq0\)；它不保证逆 Fourier 核处处非负。对复值径向权，只要以 \(|w|\) 定义的 \(I_{|w|}\) 有限，同样把配对中的相位因子放在两个平方因子之间，仍得到第二项上界。

\Needspace{6\baselineskip}
\begin{corollary}[薄球壳的两种上界]\label{b2:cor:24-thin-shell}
设整数 \(d\geq2\)，\(p_*=2(d+1)/(d+3)\)、\(q_*=p_*'\)，且 \(0<\varepsilon<1/2\)。记
\[
 \Omega_\varepsilon
 =\{\,\xi\in\mathbb R^d:1-\varepsilon<|\xi|<1+\varepsilon\,\},\quad
 \Pi_\varepsilon f
 =\mathcal F^{-1}\!\left[\mathbf1_{\Omega_\varepsilon}\cdot\widehat f\right]
\]
先定义于 Schwartz 类。则
\[
 \|\mathbf1_{\Omega_\varepsilon}\cdot\widehat f\|_2
 \leq C_d\varepsilon^{1/2}\|f\|_{p_*},\quad
 \|\Pi_\varepsilon f\|_{q_*}
 \leq C_d\varepsilon\|f\|_{p_*}.
\]
两项分别按上述范数具有唯一连续延拓。
\end{corollary}
\begin{proof}
在\cref{b2:prop:24-radial-weight}中取
\(w=\mathbf1_{(1-\varepsilon,1+\varepsilon)}\)。
由于 \(1/2<r<3/2\)，有
\[
 I_w=\int_{1-\varepsilon}^{1+\varepsilon}r^{(d-1)/(d+1)}\,\mathrm dr
 \leq C_d\varepsilon.
\]
分别代入频率平方界与混合指数界，即得两项估计及延拓。
\end{proof}

第一式中的 \(\varepsilon^{1/2}\) 来自径向积分后开平方；第二式中的 \(\varepsilon\) 来自限制与伴随各提供一次平方根。它们控制的输出空间不同，不能仅因第二项也是频率投影，就把它读成 \(L^{p_*}\) 上的投影范数。

\ProseParagraphBreak
更一般地，若 \(w\geq0\) 支于固定半径区间 \([1/2,3/2]\)，则混合算子范数至多为
\(C_d\int w(r)\,\mathrm dr\)。在离球边界约 \(2^{-j}\) 的壳层上，\((1-r^2)^\delta\) 大小约为 \(2^{-j\delta}\)，壳层厚度约为 \(2^{-j}\)，因此该层的混合指数贡献至多为
\(C_{d,\delta}2^{-j(\delta+1)}\)。只要 \(\delta>-1\)，这些边界层就可求和。要把它用于整个球，还必须处理半径趋零的部分。

\Needspace{6\baselineskip}
\begin{proposition}[全平均的混合指数界]\label{b2:prop:24-mixed-bound}
设整数 \(d\geq2\)，\(p_*=2(d+1)/(d+3)\)、\(q_*=p_*'\)，\(\delta>-1\)、\(R>0\)。对复值 Schwartz 函数，定义
\[
 B_R^\delta f(x)
 =\int_{|\xi|<R}e^{2\pi\mathrm i x\cdot\xi}
       (1-|\xi|^2/R^2)^\delta
       \cdot\widehat f(\xi)\,\mathrm d\xi.
\]
则它有唯一连续延拓
\[
 B_R^\delta:L^{p_*}(\mathbb R^d)\longrightarrow L^{q_*}(\mathbb R^d),
\]
并满足
\[
 \|B_R^\delta f\|_{q_*}
 \leq\frac{A_d^2}{2}
       \mathrm B\!\left(\frac d{d+1},\delta+1\right)
       R^{2d/(d+1)}\|f\|_{p_*}.
\]
其中 \(A_d\) 为单位球面面积测度的 \(L^{p_*}(\mathbb R^d)\to L^2(S^{d-1})\) 限制常数，\(\mathrm B\) 为 beta 函数。
\end{proposition}
\begin{proof}
先令 \(R=1\)，在径向权命题中令 \(w(r)=(1-r^2)^\delta\) 对 \(0<r<1\) 成立，在其余半径取零。
相应的积分为
\[
 I_w=\int_0^1(1-r^2)^\delta r^{(d-1)/(d+1)}\,\mathrm dr
 =\frac12\,\mathrm B\!\left(\frac d{d+1},\delta+1\right).
\]
最后一步令 \(t=r^2\)。零端的参数 \(d/(d+1)>0\)，球面边界端的参数 \(\delta+1>0\)，所以两端均可积。\(T_r\) 的范数在 \(r\downarrow0\) 时按 \(r^{-2d/q_*}\) 增大，而体积权 \(r^{d-1}\) 将总幂次变成 \((d-1)/(d+1)>-1\)，保证了零端的可积性。

对任意 \(R>0\)，直接在径向权中取
\[
 w_R(r)=
 \begin{cases}
  (1-r^2/R^2)^\delta,&0<r<R,\\
  0,&r\geq R.
 \end{cases}
\]
换元 \(r=Rs\) 后，
\[
 I_{w_R}=R^{\,1+(d-1)/(d+1)}I_w
        =R^{2d/(d+1)}I_w.
\]
再次应用径向权命题即得所述界。测试函数稠密性给出唯一延拓。
\end{proof}

对 \(\delta\geq0\)，这里的算子与本章原定义一致；负参数只在本命题明确规定的积分与连续延拓意义下使用。若 \(-1<\delta<0\)，其符号在球面附近本性无界，因而不是有界的 \(L^2\) 同指数乘子。这与命题不矛盾，因为输入、输出空间已经改变。

\ProseParagraphBreak
半径因子也能用一般缩放核对：
\[
 d\left(\frac1{p_*}-\frac1{q_*}\right)=\frac{2d}{d+1}.
\]
同指数时两个范数的尺度因子相消；这里并不相消。对单位半径证明的混合估计，不能在推广到任意半径时删掉这个幂次。

\subsection{乘子估计}
\label{b2:subsec:240403}

在无限测度空间 \(\mathbb R^d\) 上，\(L^{q_*}\) 不包含 \(L^{p_*}\)。例如选
\[
 \frac d{q_*}<a\leq\frac d{p_*},
\]
则 \((1+|x|)^{-a}\) 属于 \(L^{q_*}\)，却不属于 \(L^{p_*}\)。因而前面的混合界没有一个免费的输出空间嵌入，可以把它改成同指数界。

\ProseParagraphBreak
把输出局部化到空间球，可以明确看见缺失的常数。对半径 \(L>0\) 的球 \(B(0,L)\)，Hölder 给出
\[
 \|F\|_{L^{p_*}(B(0,L))}
 \leq |B(0,L)|^{1/p_*-1/q_*}
       \|F\|_{L^{q_*}(B(0,L))}
 =c_d L^{2d/(d+1)}\|F\|_{L^{q_*}(B(0,L))}.
\]
随着空间球增大，这个体积因子也增大。薄壳投影因此只能直接给出
\[
 \|\Pi_\varepsilon f\|_{L^{p_*}(B(0,L))}
 \leq C_d\varepsilon L^{2d/(d+1)}\|f\|_{p_*}.
\]
例如取 \(L=\varepsilon^{-1}\)，右侧的系数是
\(\varepsilon^{-(d-1)/(d+1)}\)，并未随壳厚趋零而减小。这只是球内估计；球外输出尚未处理，不能把它宣布为全空间乘子界。

\ProseParagraphBreak
有限个壳层也能呈现同一问题。设
\[
 \Lambda_j
 =\{\,\xi:1-2^{-j}<|\xi|<1-2^{-j-1}\,\},\quad
 Q_jf=\mathcal F^{-1}(\mathbf1_{\Lambda_j}\cdot\widehat f),
 \quad j\geq1.
\]
由径向权估计，对任意有限 \(N\) 和复系数 \(a_1,\ldots,a_N\)，有
\[
 \left\|\sum_{j=1}^N a_jQ_jf\right\|_{q_*}
 \leq C_d\left(\sum_{j=1}^N|a_j|2^{-j}\right)\|f\|_{p_*}.
\]
即使 \(a_j=1\)，右侧也对 \(N\) 一致；但换成空间球上的 \(p_*\) 范数时，仍须乘 \(L^{2d/(d+1)}\)。当 \(a_j\) 的大小为 \(2^{-j\delta}\) 时，混合范数中的可和条件是 \(\delta>-1\)。这个求和没有提供足以取消体积因子的额外信息。

\ProseParagraphBreak
频率局限于同一固定薄壳也不能恢复所缺的嵌入。固定 \(0<\varepsilon<1/2\)，取非零
\(\widehat\phi\in C_c^\infty(\Omega_\varepsilon)\)，于是 \(\phi\in\mathcal S\) 且 \(\Pi_\varepsilon\phi=\phi\)。对正整数 \(N\) 与距离参数 \(L\)，令
\[
 F_{N,L}(x)=\sum_{k=1}^N\phi(x-kLe_1).
\]
它的 Fourier 支集仍在同一薄壳内。固定 \(N\)，当 \(L\to\infty\) 时，对 \(s=p_*\) 与 \(s=q_*\) 分别有
\[
 \|F_{N,L}\|_s\longrightarrow N^{1/s}\|\phi\|_s.
\]
证明只需先把 \(\phi\) 在 \(L^s\) 中逼近为紧支撑函数；距离充分大时，各平移的紧支撑互不相交，范数的 \(s\) 次方便准确相加。原函数与近似函数的和相差至多 \(N\) 倍单项误差；对固定 \(N\) 再令误差趋零，即得极限。

\ProseParagraphBreak
所以在这些薄壳频率函数上，
\[
 \frac{\|F_{N,L}\|_{p_*}}{\|F_{N,L}\|_{q_*}}
 \longrightarrow
 N^{1/p_*-1/q_*}\frac{\|\phi\|_{p_*}}{\|\phi\|_{q_*}}
 =N^{2/(d+1)}\frac{\|\phi\|_{p_*}}{\|\phi\|_{q_*}}.
\]
先取大 \(N\)，再取足够大的有限距离，就能使比值任意大。这是一个有限平移模型，说明即使只考察同一频率壳层内的输出，也没有统一的 \(L^{q_*}\to L^{p_*}\) 范数转换。它不是薄壳投影同指数不有界性的独立证明；它排除的是从已有混合估计直接转换指数的那一步。

\ProseParagraphBreak
要研究同指数 \(L^p\) 界，通常还需按空间尺度分解核，把限制估计与核尾部估计结合，再控制各壳层的组合。\chapref{b2:ch:19}的 Littlewood--Paley 分解把以原点为中心、大小按倍数变化的频率环带组织起来；这里的各层都贴近同一单位球面，径向宽度不断缩小，不能不加说明地套用同一平方函数结论。不交频率在 Plancherel 中确实正交，但这只给 \(p=2\) 的平方范数，其他指数还要处理方向与空间位置上的相干。

\subsection{曲率}
\label{b2:subsec:240404}

球、椭球和多面体有助于分清局部几何与全局组合。先看线性换元：它改变球面在欧氏坐标中的形状，却不改变同指数乘子范数。

\Needspace{6\baselineskip}
\begin{proposition}[可逆线性换元下的乘子共轭]\label{b2:prop:24-linear-conjugacy}
设整数 \(d\geq1\)，\(A\) 为实可逆 \(d\times d\) 矩阵，\(m\in L^\infty(\mathbb R^d)\)，并令
\[
 T_mf=\mathcal F^{-1}(m\cdot\widehat f),\quad
 U_Af(x)=f(A^{\mathsf T}x).
\]
若 \(1<p,q<\infty\)，则 \(T_m:L^p\to L^q\) 有界，当且仅当 \(T_{m\circ A}:L^p\to L^q\) 有界；在有界情形下，
\[
 T_{m\circ A}=U_A^{-1}T_mU_A,\quad
 \|T_{m\circ A}\|_{p\to q}
 =|\det A|^{1/q-1/p}\cdot\|T_m\|_{p\to q}.
\]
这里算子均指从 Schwartz 类的连续延拓。
\end{proposition}
\begin{proof}
Fourier 换元公式与空间范数换元分别为
\[
 \widehat{U_Af}(\xi)=|\det A|^{-1}\cdot\widehat f(A^{-1}\xi),\quad
 \|U_Af\|_s=|\det A|^{-1/s}\cdot\|f\|_s
 \quad(1\leq s<\infty).
\]
因此 \(U_A^{-1}T_mU_A\) 的符号为 \(m(A\xi)\)，得到共轭恒等式。若 \(T_m\) 有界，两个拉回范数给出
\[
 \|T_{m\circ A}\|_{p\to q}
 \leq|\det A|^{1/q-1/p}\cdot\|T_m\|_{p\to q}.
\]
反向对 \(A^{-1}\) 和符号 \(m\circ A\) 应用同一结论，得到逆向不等式及有界性的等价。测试函数稠密性使共轭关系保持到两个 Banach 空间之间。
\end{proof}

取 \(m(\xi)=(1-|\xi|^2)_+^\delta\)、\(\delta\geq0\)，就得到椭球符号
\((1-|A\xi|^2)_+^\delta\) 的平均。在同指数 \(p=q\) 时，行列式因子完全相消；因此固定椭球与球具有相同的有界性范围和算子范数。在不同指数上则保留这个因子，与前面球半径的缩放一致。若另外估计各欧氏方向导数，矩阵与其逆矩阵会进入导数公式；这种正则性常数的依赖不能与纯 \(L^p\) 乘子范数混为一谈。

\ProseParagraphBreak
平面边界提供另一种可直接构造的模型。需要恢复的前章工具是一维 Hilbert 变换：
\[
 \widehat{Hf}(\xi)=-\mathrm i\,\operatorname{sgn}(\xi)
                  \cdot\widehat f(\xi),\quad
 \|Hf\|_{L^p(\mathbb R)}\leq C_{H,p}\|f\|_{L^p(\mathbb R)}
 \quad(1<p<\infty).
\]
其平方定义见\cref{b2:thm:18-hilbert-l2}；核 \(1/(\pi x)\) 的大小与差分满足 Calderón--Zygmund 条件，所以\cref{b2:thm:18-cz-lp}给出中间指数有界性及与平方定义的兼容。由于 \(\mathrm iH\) 的符号为 \(\operatorname{sgn}\xi\)，负半轴投影为 \((I-\mathrm iH)/2\)。在高维中沿固定方向应用这个一维算子，Fubini 保持同一个 \(L^p\) 上界。

\Needspace{6\baselineskip}
\begin{proposition}[有限半空间交的投影]\label{b2:prop:24-polytope-projection}
设整数 \(d\geq1\)，\(1<p<\infty\)，正整数 \(N\)，单位向量
\(\nu_1,\ldots,\nu_N\in S^{d-1}\) 及实数 \(b_1,\ldots,b_N\) 给出
\[
 P=\bigcap_{j=1}^N\{\,\xi\in\mathbb R^d:\xi\cdot\nu_j<b_j\,\}.
\]
令 \(T_Pf=\mathcal F^{-1}(\mathbf1_P\cdot\widehat f)\) 先在 Schwartz 类上定义。则 \(T_P\) 唯一连续延拓为 \(L^p(\mathbb R^d)\) 上的有界算子，并且
\[
 \|T_P\|_{p\to p}\leq
 \left(\frac{1+C_{H,p}}2\right)^N.
\]
这里 \(C_{H,p}\) 是一维 Hilbert 变换的 \(L^p\) 范数上界；常数不依赖法向与各平面的位置。集合 \(P\) 可以无界，空集情形也包括在内。
\end{proposition}
\begin{proof}
对单位向量 \(\nu\)，用旋转把一维 Hilbert 变换放到该方向，记作 \(H_\nu\)。它在平方空间的符号为
\(-\mathrm i\,\operatorname{sgn}(\xi\cdot\nu)\)，且 Fubini 给出
\(\|H_\nu\|_{p\to p}\leq C_{H,p}\)。于是
\[
 Q_{\nu,0}=\frac{I-\mathrm iH_\nu}{2}
\]
在 \(L^2\) 上是 \(\{\xi\cdot\nu<0\}\) 的投影，在 \(L^p\) 上有界至多 \((1+C_{H,p})/2\)。令
\[
 M_{b\nu}f(x)=e^{2\pi\mathrm i b\nu\cdot x}\cdot f(x).
\]
调制在 \(L^p\) 上等距，且
\(Q_{\nu,b}=M_{b\nu}Q_{\nu,0}M_{-b\nu}\)
的符号为 \(\mathbf1_{\{\xi\cdot\nu<b\}}\)。

在 \(L^2\) 上，这些算子都是 Fourier 乘法，因此彼此交换，有限复合的符号为
\[
 \prod_{j=1}^N\mathbf1_{\{\xi\cdot\nu_j<b_j\}}=\mathbf1_P.
\]
每项的 \(L^p\) 延拓与其 \(L^2\) 定义在交集上一致，所以对 Schwartz 输入，复合既等于 \(T_P\)，又有上述 \(N\) 个常数的乘积界。由 \(L^p\) 稠密性完成延拓。边界包含在有限个超平面中，Lebesgue 测度为零，故使用开或闭半空间不改变乘子。
\end{proof}

矩形是这个结论的一个特例，也可以直接逐坐标使用区间投影。一般有限凸多面体的各个平面方向可以不同，但有限次复合仍有界。这个证明给出的是随面数积累的上界，尚未说明它在越来越多面时是否真的必须变大。下面用逼近球的有限模型回答后一问题。

\Needspace{6\baselineskip}
\begin{proposition}[逼近球的多面体范数发散]\label{b2:prop:24-polytope-divergence}
设整数 \(d\geq2\)，\(1<p<\infty\)、\(p\neq2\)。取 \(0<\eta_k<1\)、\(\eta_k\to0\)，并取单位球面的有限 \(\eta_k\)-网 \(V_k\subset S^{d-1}\)，即每个单位向量到 \(V_k\) 某点的欧氏距离小于 \(\eta_k\)。定义外切凸多面体
\[
 P_k=\bigcap_{\nu\in V_k}\{\,\xi:\xi\cdot\nu<1\,\},\quad
 T_{P_k}f=\mathcal F^{-1}(\mathbf1_{P_k}\cdot\widehat f).
\]
则每个 \(T_{P_k}\) 都在 \(L^p(\mathbb R^d)\) 上有界，但
\[
 \|T_{P_k}\|_{p\to p}\longrightarrow\infty.
\]
\end{proposition}
\begin{proof}
有限网由单位球面的紧致性取得。单位开球包含于 \(P_k\)。若 \(\xi\in P_k\)、\(\xi\neq0\)，令 \(u=\xi/|\xi|\)，并选 \(\nu\in V_k\) 满足 \(|u-\nu|<\eta_k\)。单位向量的恒等式给出
\[
 u\cdot\nu=1-\frac{|u-\nu|^2}{2}
 >1-\frac{\eta_k^2}{2}.
\]
结合 \(\xi\cdot\nu<1\)，得到
\[
 B(0,1)\subset P_k
 \subset B\!\left(0,(1-\eta_k^2/2)^{-1}\right).
\]
所以各 \(P_k\) 有界，且其示性函数在单位球面之外逐点趋于 \(\mathbf1_{B(0,1)}\)。球面是 Lebesgue 零测集，故符号几乎处处收敛。

每个多面体投影的有界性已由\cref{b2:prop:24-polytope-projection}证明。若所述范数不趋于无穷，则存在一个子列，其范数统一不超过某个有限 \(M\)。对任意固定 \(f\in\mathcal S\)，符号差的模至多为一，且 \(\widehat f\in L^1\)，故支配收敛给出
\[
 \|( \mathbf1_{P_k}-\mathbf1_{B(0,1)})\cdot\widehat f\|_1\longrightarrow0.
\]
因而这一子列的 \(T_{P_k}f\) 一致收敛到硬球投影 \(B_1^0f\)。Fatou 引理于是给出
\[
 \|B_1^0f\|_p^p
 \leq\liminf_k\|T_{P_k}f\|_p^p
 \leq M^p\|f\|_p^p.
\]
稠密性将硬球投影延拓为有界 \(L^p\) 算子，与\cref{b2:thm:24-fefferman}在 \(d\geq2\)、\(p\neq2\) 时的结论矛盾。
\end{proof}

这个发散结论使用了本章明确标出的深层球乘子定理，但由它到有限多面体模型的极限步骤都已写出。它没有提供关于面数的定量发散速度；在 \(p=2\) 时，各个非空多面体投影的范数反而恒为一。几何区域逐点逼近球，并不能提供非二指数上一致的算子常数。

\ProseParagraphBreak
曲率的两种作用因此需要分开理解。对一张固定的驻相片，非零曲率能改善振荡核衰减；当越来越多方向共同出现时，输出又可能形成多方向集中。前一种效应服务于正平滑平均的核估计，后一种效应参与硬球投影的失败。分析新的频率区域时，应同时考察单片衰减、片的数量以及各片怎样组合；只有局部衰减还不够。

\ProseParagraphBreak
本章各项结果也各有明确的适用范围：尾部检验给出必要条件，可积核与解析插值给出同指数充分范围，球面限制给出本节的混合指数界，球乘子定理排除高维硬截断的非二指数。应用文献中的进一步估计时，还应核对输入是否径向、空间是否局部、半径是否统一及是否允许尺度损失，不能把改变了其中条件的结论直接填入普通全空间 \(L^p\) 乘子问题。


\begin{chaptersummary}
Bochner--Riesz 符号的球边界决定核的远场振荡，原点附近的展开则决定光滑输入的低频误差。前者控制可积性与乘子障碍，后者可以通过半径组合改变。两种机制虽然出现在同一个符号中，却需要不同的估计；更高的边界消失阶并不无条件意味着更快的逼近。

\ProseParagraphBreak

可积核给出全部普通指数的稳定性、有限指数下的范数逼近以及由极大函数控制的逐点结论。临界核失去绝对可积性以后，这套证明停止，但平方理论仍然成立。尾部检验、解析插值与高维硬球否定定理各自提供不同信息，留下的正阶区间不能靠形式代入或指数包含补齐。

\ProseParagraphBreak

限制理论通过壳层给出混合指数估计。将它转成同指数乘子界，必须另外处理空间局部化、尺度求和和方向组合。固定多面体的有界性也不能在缺少一致常数时传递给球。无论改变半径、阶数还是几何形状，都应同时追踪算子常数和收敛方式，才能判断一个极限真正保留了什么。
\end{chaptersummary}

\BookExercises

% 自拟；中等：固定频带上的阶数半群、生成元定义域与准确范数差，区分子空间恒等和全空间投影。
\Needspace{6\baselineskip}
\begin{exercise}\label{b2:ex:24-new-1}
设整数 \(d\geq1\)、\(R>0\)，令
\[
 \mathcal H_R=\{\,f\in L^2(\mathbb R^d):
          \widehat f=0\ \text{几乎处处于 }\{|\xi|\geq R\}\,\}.
\]
把 \(B_R^\delta\) 限制到这个 Hilbert 子空间，证明 \(\delta\geq0\) 构成强连续的压缩半群。若生成元按
\(Gf=\lim_{h\downarrow0}(B_R^hf-f)/h\) 定义，证明
\[
 \widehat{Gf}(\xi)=\log(1-|\xi|^2/R^2)\cdot\widehat f(\xi)
       \quad(|\xi|<R),
\]
球外取零；其定义域恰为使右侧在球内属于 \(L^2\) 的 \(f\in\mathcal H_R\)。再证明对 \(0<\delta<\eta\)，
\[
 \|B_R^\delta-B_R^\eta\|_{\mathcal H_R\to\mathcal H_R}
 =\left(1-\frac\delta\eta\right)
       \left(\frac\delta\eta\right)^{\delta/(\eta-\delta)},
\]
而 \(\|B_R^\delta-I\|_{\mathcal H_R\to\mathcal H_R}=1\) 对每个 \(\delta>0\) 成立。说明为什么同一个族在全空间上的零阶算子不能被称为恒等。
\end{exercise}
\begin{solution}
球内令 \(a(\xi)=1-|\xi|^2/R^2\)，则 \(0<a\leq1\)。符号相乘给出
\(B_R^\delta B_R^\eta=B_R^{\delta+\eta}\)，其模不超过一；零阶符号在 \(\mathcal H_R\) 上为一。对每个固定输入，\(a^h\to1\) 且 \(|a^h-1|\leq1\)，Plancherel 与支配收敛证明零点强连续。半群性质和压缩界把连续性移到其余参数。

记 \(\ell(\xi)=-\log a(\xi)\geq0\)。当 \(\ell\widehat f\in L^2\) 时，
\[
 \frac{e^{-h\ell}-1}{h}\widehat f\longrightarrow-\ell\widehat f,
 \quad
 \left|\frac{e^{-h\ell}-1}{h}\widehat f\right|
       \leq\ell|\widehat f|.
\]
平方支配收敛给出生成元极限。反过来，若差商在 \(L^2\) 中有极限，取一个几乎处处收敛的子列；其点态极限只能是 \(-\ell\widehat f\)，所以这个函数必在 \(L^2\)，也必是原来的平方极限。定义域的两向包含均已得到。

算子范数由符号的本性上确界给出。函数 \(u^\delta-u^\eta\) 在 \(0\leq u\leq1\) 非负，内部最大点满足
\(u^{\eta-\delta}=\delta/\eta\)。代入便得题中准确公式；相应的 \(u\) 值对应正测度邻域的球内频率，所以也是本性上确界。当比较正阶与零阶时，\(1-u^\delta\) 在 \(u\downarrow0\) 趋于一，故范数差恒为一。这不妨碍每个固定输入强连续。

在整个 \(L^2\) 上，\(B_R^0\) 是频率球的正交投影，消去了球外数据，因而不是恒等。只有先明确限制到 \(\mathcal H_R\)，上述族才具有半群在零参数处所需的恒等元。
\end{solution}

% 自拟；中等：半径趋零时的质量扩散轮廓，求L1极限范数并与中间指数的强趋零比较。
\Needspace{6\baselineskip}
\begin{exercise}\label{b2:ex:24-new-2}
设整数 \(d\geq1\)、\(\delta>(d-1)/2\)，记 \(K_\delta\) 为单位半径核。对 \(f\in L^1(\mathbb R^d)\)，令 \(M_f=\int f\)。证明当 \(R\downarrow0\) 时，对每个 \(1\leq q\leq\infty\)，
\[
 R^{-d}(B_R^\delta f)(\,\cdot\,/R)
       \longrightarrow M_fK_\delta\quad\text{于 }L^q,
\]
其中 \(q=\infty\) 指一致范数。由此证明
\[
 \lim_{R\downarrow0}R^{-d(1-1/q)}\|B_R^\delta f\|_q
       =|M_f|\cdot\|K_\delta\|_q.
\]
说明 \(L^1\) 输入的输出何时强趋于零，并证明对每个 \(1<p<\infty\) 及每个 \(f\in L^p\)，都有 \(B_R^\delta f\to0\) 于 \(L^p\)。
\end{exercise}
\begin{solution}
可积核定理和频率符号可积性说明 \(K_\delta\in L^1\cap C_0\)，因而属于全部 \(L^q\)、\(1\leq q\leq\infty\)。它的平移在有限 \(q\) 范数中连续；在一致范数中也连续，因为 \(C_0\) 函数一致连续。缩放核公式给出
\[
 R^{-d}(B_R^\delta f)(z/R)=\int K_\delta(z-Ry)\cdot f(y)\,\mathrm dy.
\]
减去 \(M_fK_\delta(z)\)，再用 Minkowski 或一致范数的积分上界，得到
\[
 \left\|R^{-d}(B_R^\delta f)(\,\cdot\,/R)-M_fK_\delta\right\|_q
 \leq\int|f(y)|\cdot
          \|K_\delta(\,\cdot-Ry)-K_\delta\|_q\,\mathrm dy.
\]
对固定 \(y\)，范数趋于零；它至多为 \(2\|K_\delta\|_q\)。以 \(|f|\) 为可积控制，支配收敛就给出第一条结论。空间换元把左侧未减去轮廓的范数写成
\(R^{-d(1-1/q)}\|B_R^\delta f\|_q\)，故也得到第二条结论。

特别地，\(\|B_R^\delta f\|_1\to|M_f|\cdot\|K_\delta\|_1\)，所以强趋于零恰好要求 \(M_f=0\)。当 \(M_f\ne0\) 时甚至不存在任何 \(L^1\) 强极限：一方面
\(\|B_R^\delta f\|_\infty\leq R^d\|K_\delta\|_\infty\cdot\|f\|_1\to0\)，
另一方面其 \(L^1\) 范数不趋零。若有强极限，几乎处处收敛子列会迫使极限为零，形成矛盾。

若 \(1<p<\infty\) 且 \(f\in L^1\cap L^p\)，上面第二式的 \(q=p\) 情形与 \(d(1-1/p)>0\) 给出 \(L^p\) 强趋零。一般 \(L^p\) 输入可用 \(L^1\cap L^p\) 逼近，再用 \(\sup_R\|B_R^\delta\|_{p\to p}\leq\|K_\delta\|_1\) 推广。这里的小半径把质量摊到越来越大的空间尺度；总质量仍可保留，而中间指数范数已经趋零。
\end{solution}

% 自拟；综合计算：热模型的L²算子范数收敛及一阶修正，用一致尾部控制升级点态Taylor展开。
\Needspace{6\baselineskip}
\begin{exercise}\label{b2:ex:24-new-3}
设整数 \(d\geq1\)，\(H\) 为符号 \(e^{-|\xi|^2}\) 的热算子，\(Q\) 为符号 \(\tfrac12|\xi|^4e^{-|\xi|^2}\) 的平方乘子。证明
\[
 \left\|n\bigl(B_{\sqrt n}^n-H\bigr)+Q\right\|_{2\to2}
       \longrightarrow0\quad(n\to\infty),
\]
并求出准确的一阶范数极限
\[
 \lim_{n\to\infty}n\|B_{\sqrt n}^n-H\|_{2\to2}=2e^{-2}.
\]
需要证明对全部频率的一致控制，仅在固定紧频带内展开并不足够。
\end{exercise}
\begin{solution}
令 \(a=|\xi|^2\)，并记
\[
 D_n(a)=n\left[e^{-a}-(1-a/n)_+^n\right]\geq0.
\]
在任意固定的紧 \(a\) 区间内，Taylor 展开一致给出
\[
 n\log(1-a/n)=-a-\frac{a^2}{2n}+O(n^{-2}),
 \quad D_n(a)\longrightarrow\frac{a^2}2e^{-a}.
\]
余项常数在这一步允许依赖所固定的紧区间。

为控制其外部，先在 \(0\leq a\leq n/2\) 置
\(u_n=-n\log(1-a/n)-a\)。对数级数给出
\[
 0\leq u_n
 =n\sum_{k=2}^\infty\frac{(a/n)^k}{k}
 \leq\frac{a^2}{2n(1-a/n)}\leq\frac{a^2}n.
\]
所以
\(D_n(a)=ne^{-a}(1-e^{-u_n})\leq a^2e^{-a}\)。若 \(a>n/2\)，则直接有
\(D_n(a)\leq ne^{-a}<2ae^{-a}\)，包括 \(a\geq n\) 的截断区域。因此整个半轴上的尾部都由
\((a^2+2a)e^{-a}\) 一致控制，而这个函数在无穷远趋零。先选大紧区间使尾部小，再在紧区间内取极限，即得
\[
 \sup_{a\geq0}\left|D_n(a)-\frac{a^2}2e^{-a}\right|\longrightarrow0.
\]
Plancherel 将符号的一致收敛转为第一条算子范数收敛。最后，\(a^2e^{-a}/2\) 在 \(a=2\) 处取得最大值 \(2e^{-2}\)；由一致收敛，其上确界也是 \(D_n\) 上确界的极限，得到所求准确常数。
\end{solution}

% 自拟；反例训练：临界核的弱L1性质不自动给卷积弱(1,1)，以正核和扩张输入展示端点Young失效。
\Needspace{6\baselineskip}
\begin{exercise}\label{b2:ex:24-new-4}
设整数 \(d\geq1\)、\(\delta_c=(d-1)/2\)。对可测函数 \(u\)，记
\[
 \|u\|_{L^{1,\infty}}
 =\sup_{\lambda>0}\lambda\cdot
       |\{\,x:|u(x)|>\lambda\,\}|.
\]
证明临界 Bochner--Riesz 核满足
\(K_{\delta_c}\in L^{1,\infty}\setminus L^1\)。再考虑正核
\(k(x)=(1+|x|)^{-d}\)，证明它也属于弱 \(L^1\)，但不存在常数 \(C\) 使
\[
 \|k*f\|_{L^{1,\infty}}\leq C\|f\|_1
\]
对所有非负光滑紧支撑输入成立。说明这个反例对临界 Bochner--Riesz 算子本身能够、不能够推出什么。
\end{exercise}
\begin{solution}
临界核渐近和原点处的有界性给出
\(|K_{\delta_c}(x)|\leq C(1+|x|)^{-d}\)。因此其超过 \(\lambda\) 的集合包含于一个体积至多为 \(C'/\lambda\) 的球，弱 \(L^1\) 准范数有限。另一方面，振荡主项在每个充分大的单位径向区间中都有固定比例的非零下界，径向绝对积分至少贡献 \(c/r\)；调和级数发散，所以它不在 \(L^1\)。

同一个分布函数上界说明 \(k\in L^{1,\infty}\)。取非负 \(\psi\in C_c^\infty(\mathbb R^d)\)，使 \(\int\psi=1\)，且在单位球上 \(\psi\geq c_0>0\)，令
\[
 f_R(x)=R^{-d}\cdot\psi(x/R),\quad R\geq4.
\]
这些输入的 \(L^1\) 范数为一。若 \(|x|<R/2\)，在 \(|y|<R\) 上 \(f_R(y)\geq c_0R^{-d}\)，而平移后的积分区域包含 \(B(0,R/2)\)，故
\[
 (k*f_R)(x)
 \geq c_0R^{-d}\int_{|z|<R/2}(1+|z|)^{-d}\,\mathrm dz
 \geq c_dR^{-d}\log R.
\]
最后一个下界来自 \(1\leq|z|\leq R/2\) 上与 \(\mathrm dr/r\) 可比的径向积分。取 \(\lambda\) 为右侧的一半，超水平集的测度至少为 \(|B(0,R/2)|\)，所以
\[
 \|k*f_R\|_{L^{1,\infty}}\geq c'_d\log R\longrightarrow\infty.
\]
这排除了所声称的弱型界。

反例说明不能把 Young 的 \(L^1\) 核假设机械改成弱 \(L^1\) 核假设，并在端点保留同样的结论。它没有证明临界 Bochner--Riesz 算子无弱型界，因为 \(k\) 是另一个没有振荡抵消的核；对原算子的端点判断仍需利用它自身的结构。
\end{solution}

% 自拟；综合证明：保留三半径消项并推广有限半径外推，证明导数界、饱和阶及整数参数例外。
\Needspace{6\baselineskip}
\begin{exercise}\label{b2:ex:24-new-5}
设整数 \(d\geq1\)、\(1<p<\infty\)、\(\delta>0\)，并假定
\(\sup_{R>0}\|B_R^\delta\|_{p\to p}<\infty\)。先求一组不依赖 \(\delta\) 的三个权，使
\[
 C_R=w_1B_R^\delta+w_2B_{\sqrt2R}^\delta+w_3B_{\sqrt3R}^\delta
\]
的低频符号对每个 \(\delta>0\) 都等于 \(1+O(|\xi|^6/R^6)\)，并证明满足这一统一要求的权唯一。进一步，对整数 \(m\geq2\)，证明
\[
 w_j=\frac{(-1)^{m-j}j^{m-1}}{(j-1)!(m-j)!},
 \quad C_R^{(m)}=\sum_{j=1}^m w_j B_{\sqrt jR}^\delta
\]
满足
\[
 \|C_R^{(m)}f-f\|_p\leq C R^{-s}\|f\|_{H^{s,p}},
 \quad R\geq1,\quad0<s\leq2m.
\]
这里 \(H^{s,p}\) 使用 \(J^s=(I-\Delta)^{s/2}\) 的范数；常数可依赖 \(d,p,\delta,m,s\) 和已有的一致算子界，不依赖 \(R,f\)。求第一个可能非零的 Taylor 系数，证明当 \(\delta\notin\{1,\ldots,m-1\}\) 时通常不能超过 \(2m\) 阶；当 \(\delta\) 属于这个有限整数集时，证明上述界反而允许任意 \(s>0\)。
\end{exercise}
\begin{solution}
令 \(v=|\xi|^2/R^2\)。要保留常数并消去 \(v,v^2\)，相应的插值权应满足
\[
 w_1+w_2+w_3=1,
 \quad w_1+\frac{w_2}2+\frac{w_3}3=0,
 \quad w_1+\frac{w_2}4+\frac{w_3}9=0.
\]
方程的唯一解为 \((1/2,-4,9/2)\)。若 \(\delta=1\)，原符号本无二次项；这里的三条条件仍选出同一组通用于所有 \(\delta\) 的插值权，单独要求这一特殊参数的误差阶则不会唯一确定三个权。

一般情形取节点 \(t_j=1/j\)。次数至多 \(m-1\) 的多项式 \(P\) 满足 Lagrange 插值恒等式
\[
 P(0)=\sum_{j=1}^mP(t_j)
          \prod_{k\ne j}\frac{-t_k}{t_j-t_k}.
\]
直接约分所得乘积，就是题中的 \(w_j\)。因此
\[
 \sum_jw_j=1,
 \quad \sum_jw_jj^{-\ell}=0\quad(1\leq\ell<m).
\]
这些条件唯一确定权重，因为节点互异。对 \(P(t)=t^m\)，其插值误差是 \(\prod_j(t-t_j)\)，在零点取值得
\[
 \sum_jw_jj^{-m}=\frac{(-1)^{m-1}}{m!}.
\]
在 \(0\leq v<1\) 内逐项展开，因而有
\[
 \sum_{j=1}^m w_j(1-v/j)^\delta
 =1-\frac{\delta(\delta-1)\cdots(\delta-m+1)}{(m!)^2}v^m
       +O(v^{m+1}).
\]
对固定 \(m,\delta\)，余项在 \(0\leq v\leq1/4\) 上连同任意有限阶导数受控。

下面证明整个 \(L^p\) 误差界。首先，组合算子的统一范数至多为
\(\sum_j|w_j|\) 乘原来的统一界。取光滑截断 \(\chi\)，在 \(|\xi|\leq1/4\) 为一，支撑于 \(|\xi|<1/2\)，并令 \(g=J^sf\)。高频输入可写成 \(R^{-s}T_{a_R}g\)，其中
\[
 a_R(\xi)=R^s\bigl(1-\chi(\xi/R)\bigr)
                         (1+4\pi^2|\xi|^2)^{-s/2}.
\]
在其非零区域有 \(|\xi|\geq R/4\)。对每个多重指标 \(\alpha\)，乘积与链式法则给出
\(|\xi|^{|\alpha|}|\partial^\alpha a_R(\xi)|\leq C_{\alpha,s}\)：
落在截断上的导数只出现于 \(|\xi|\asymp R\)，落在 Bessel 权上的导数则每次提供一个由 \(|\xi|^{-1}\) 控制的因子。故 Mihlin 定理给出一致界，再应用 \(C_R^{(m)}-I\) 的一致界，得到高频误差。

低频误差符号可以写成 \(v^m b(v)\)，其中 \(b\) 在上述紧区间光滑。除去所求的 \(R^{-s}\) 后，对 \(g\) 的符号为
\[
 b_R(\xi)=R^{s-2m}|\xi|^{2m}
    \cdot b(|\xi|^2/R^2)\chi(\xi/R)
                     (1+4\pi^2|\xi|^2)^{-s/2}.
\]
给出其尺度导数界可以明确看见 \(2m\) 的作用。多项式 \(|\xi|^{2m}\) 的每个非零 \(k\) 阶导数至多为 \(C_k|\xi|^{2m-k}\)；归一化光滑因子的导数至多为 \(C_\alpha R^{-|\alpha|}\)，在 \(0<|\xi|\leq R/2\) 上可进一步用 \(C_\alpha|\xi|^{-|\alpha|}\) 控制。Bessel 权的导数也满足相应的尺度界。于是乘积求导后有
\[
 |\xi|^{|\alpha|}\cdot|\partial^\alpha b_R(\xi)|
 \leq C_{\alpha,m,\delta,s}
        R^{s-2m}|\xi|^{2m}(1+|\xi|)^{-s}.
\]
当 \(|\xi|\leq1\) 时，右侧由 \(CR^{s-2m}\) 控制；当 \(1\leq|\xi|\leq R/2\) 时，由 \(C(|\xi|/R)^{2m-s}\) 控制。若 \(s\leq2m\)、\(R\geq1\)，两者一致有界。取到 Mihlin 定理需要的阶数，低频误差也得到 \(CR^{-s}\|g\|_p\) 上界，完成证明。

若 \(\delta\notin\{1,\ldots,m-1\}\)，上述首项系数非零。固定一个非零 Schwartz 输入，使 \(\widehat f\) 光滑紧支撑于远离零点的环带。固定频带上的展开连同全部有限阶导数一致成立，所以逆变换甚至在 Schwartz 拓扑中收敛，并有
\[
 R^{2m}(C_R^{(m)}-I)f\longrightarrow
 -\frac{\delta(\delta-1)\cdots(\delta-m+1)}
          {(m!)^2(2\pi)^{2m}}\,(-\Delta)^mf.
\]
右端非零，排除了对所有这类光滑输入都成立的更高幂次速度。若 \(\delta\in\{1,\ldots,m-1\}\)，则每个 \((1-v/j)^\delta\) 在 \(v<1\) 是次数低于 \(m\) 的多项式，插值条件使组合精确等于一。低频误差于是完全消失；前面的高频论证对每个 \(s>0\) 都适用，给出最后的例外结论。
\end{solution}

% 自拟；综合计算：可积核阈值附近端点范数的准确爆破常数，使用一致渐近和周期均值而非只判发散。
\Needspace{6\baselineskip}
\begin{exercise}\label{b2:ex:24-new-6}
设整数 \(d\geq1\)，令 \(\delta_c=(d-1)/2\)。利用正文中对紧阶数区间一致成立的核渐近，证明
\[
 \lim_{\varepsilon\downarrow0}
      \varepsilon\|K_{\delta_c+\varepsilon}\|_1
 =\frac{4\Gamma((d+1)/2)}{\pi^{3/2}\Gamma(d/2)}.
\]
因此两个端点 \(L^1\)、\(L^\infty\) 的算子范数在接近可积性阈值时都按 \(\varepsilon^{-1}\) 增长。证明中应说明为什么余弦的相位随 \(\varepsilon\) 变化不影响极限。
\end{exercise}
\begin{solution}
记 \(\omega_{d-1}=2\pi^{d/2}/\Gamma(d/2)\)，并令
\[
 a_\delta=\frac{\Gamma(\delta+1)}{\pi^{\delta+1}},
 \quad \theta_\delta=\frac{\pi(d+1+2\delta)}4.
\]
对 \(0\leq\varepsilon\leq1\)、\(r\geq1\)，一致渐近为
\[
 K_{\delta_c+\varepsilon}(r)
 =a_{\delta_c+\varepsilon}r^{-d-\varepsilon}
       \cos(2\pi r-\theta_{\delta_c+\varepsilon})
       +O(r^{-d-\varepsilon-1}).
\]
利用 \(\bigl||u+v|-|u|\bigr|\leq|v|\)，余项对径向绝对积分的贡献至多为
\(C\int_1^\infty r^{-2-\varepsilon}\,\mathrm dr\leq C\)。在单位球内，核由其频率符号的 \(L^1\) 范数一致控制，贡献也为 \(O(1)\)。所以
\[
 \|K_{\delta_c+\varepsilon}\|_1
 =\omega_{d-1}a_{\delta_c+\varepsilon}
   \int_1^\infty r^{-1-\varepsilon}
          |\cos(2\pi r-\theta_{\delta_c+\varepsilon})|\,\mathrm dr+O(1).
\]

函数 \(|\cos(2\pi r-\theta)|\) 的周期均值为 \(2/\pi\)，与 \(\theta\) 无关。减去这个均值后的函数有一个一致有界的原函数
\[
 H_\theta(r)=\int_1^r
       \left(|\cos(2\pi u-\theta)|-\frac2\pi\right)\mathrm du.
\]
有界性由整周期积分为零、剩余长度小于一直接得到，且界与相位无关。对权 \(r^{-1-\varepsilon}\) 分部积分，边界项及导数积分均一致有界，故
\[
 \int_1^\infty r^{-1-\varepsilon}|\cos(2\pi r-\theta)|\,\mathrm dr
 =\frac2{\pi\varepsilon}+O(1)
\]
对所有相位及 \(0<\varepsilon\leq1\) 一致成立。乘 \(\varepsilon\) 并令其趋零，得到极限
\[
 \frac2\pi\omega_{d-1}a_{\delta_c}
 =\frac{4\Gamma((d+1)/2)}{\pi^{3/2}\Gamma(d/2)}.
\]
最后，正文已证明可积卷积核的两个端点算子范数都恰为核的 \(L^1\) 范数；缩放不改变此范数。因此上述极限同时给出了全部半径的两个端点常数的准确增长。
\end{solution}

% 自拟；反例与泛函分析：正阶趋零时的扩展乘子范数必爆破，区分单个符号平滑、统一界及单输入失控。
\Needspace{6\baselineskip}
\begin{exercise}\label{b2:ex:24-new-7}
设整数 \(d\geq2\)、\(1<p<\infty\)、\(p\ne2\)。对 \(\delta>0\)，令
\[
 N_p(\delta)=
 \begin{cases}
 \|B_1^\delta\|_{p\to p},&B_1^\delta\text{ 有有界的 }L^p\text{ 延拓},\\
 +\infty,&\text{否则}.
 \end{cases}
\]
以正文的 Fefferman 球乘子定理为前提，证明
\(N_p(\delta)\to+\infty\) 当 \(\delta\downarrow0\)。进一步，若 \(\delta_n\downarrow0\) 是一个满足每个 \(N_p(\delta_n)<\infty\) 的序列，证明存在同一个 \(f\in L^p\)，使
\[
 \limsup_{n\to\infty}\|B_1^{\delta_n}f\|_p=+\infty.
\]
本题不假定任意正阶都已有界，也不要求构造适用于所有输入的失败结论。
\end{exercise}
\begin{solution}
若第一个结论不成立，则存在 \(\delta_n\downarrow0\) 的子序列与有限常数 \(C\)，使 \(N_p(\delta_n)\leq C\)。其符号在球内部趋于一，在球外为零，且模始终不超过一。因此对任意 Schwartz 输入，Plancherel 与支配收敛给出
\[
 B_1^{\delta_n}f\longrightarrow B_1^0f\quad\text{于 }L^2.
\]
选取几乎处处收敛的子列，再用 Fatou，得到
\[
 \|B_1^0f\|_p^p
 \leq\liminf_n\|B_1^{\delta_n}f\|_p^p
 \leq C^p\|f\|_p^p.
\]
Schwartz 类在有限指数 \(L^p\) 中稠密，因而这会给出硬球乘子的有界延拓，与\cref{b2:thm:24-fefferman}矛盾。此论证排除了所有趋零而范数有界的阶数列，正是扩展意义下的 \(N_p(\delta)\to+\infty\)。

现在假定题中全由有界算子组成的序列确实存在。若每个 \(f\in L^p\) 都满足 \(\sup_n\|B_1^{\delta_n}f\|_p<\infty\)，一致有界原理将给出
\(\sup_nN_p(\delta_n)<\infty\)，与第一步矛盾。因此存在一个输入使输出范数的上确界无穷。每个固定 \(n\) 的输出范数有限，删除任意有限前缀仍然无界，所以其上极限也为无穷。

这是一条依赖深层球乘子定理的退化结论，并没有提供爆破的定量速度。它也不证明某个给定正阶有界；若小阶数已经无界，第一条结论按定义仍然成立，第二条的序列假设则可能无法满足。量词中的统一常数和同一个输入都不能省略。
\end{solution}

% 自拟；综合应用：一维跳跃的完整边界轮廓与过冲，利用累积振荡尾部证明准确Lp误差尺度。
\Needspace{6\baselineskip}
\begin{exercise}\label{b2:ex:24-new-8}
在 \(\mathbb R\) 上固定有限实数 \(\delta>0\)，令 \(K_\delta\) 为单位半径核，并置
\[
 F_\delta(t)=\int_{-\infty}^tK_\delta(y)\,\mathrm dy,
 \quad H(t)=\mathbf1_{(0,\infty)}(t),
 \quad G_\delta=F_\delta-H.
\]
对 \(f=\mathbf1_{(0,1)}\)，证明 \(B_R^\delta f(t/R)\to F_\delta(t)\) 对每个固定 \(t\) 成立。证明 \(F_\delta\) 在某些点大于一，在某些点小于零；不能只从 \(K_\delta\) 变号推出这件事。最后证明对每个 \(1\leq p<\infty\)，
\[
 \lim_{R\to\infty}R^{1/p}\|B_R^\delta f-f\|_p
       =2^{1/p}\|G_\delta\|_p\in(0,\infty),
\]
而一致范数误差始终不小于 \(1/2\)。
\end{exercise}
\begin{solution}
一维中 \(\delta>0\) 已在可积核范围，核为偶实函数且总积分为一。因此 \(F_\delta\) 连续，满足
\[
 F_\delta(-t)=1-F_\delta(t),
 \quad F_\delta(0)=\frac12,
 \quad F_\delta(-\infty)=0,\quad F_\delta(+\infty)=1.
\]
直接对区间积分，得到
\[
 B_R^\delta f(t/R)=\int_{t-R}^tK_\delta(y)\,\mathrm dy
       \longrightarrow F_\delta(t).
\]
这给出左端跳跃在宽度 \(R^{-1}\) 内的完整轮廓。

要证明过冲，必须计算累积尾部。记
\(a_\delta=\Gamma(\delta+1)/\pi^{\delta+1}\)、
\(\theta_\delta=\pi(\delta+1)/2\)。将正文的端点展开多保留一项，得到
\[
 K_\delta(r)=a_\delta r^{-\delta-1}
 \left[\cos(2\pi r-\theta_\delta)
       -\frac{\delta(\delta+1)}{4\pi r}
           \sin(2\pi r-\theta_\delta)\right]
       +O(r^{-\delta-3}).
\]
具体地，在频率端点 \(\xi=1\) 令 \(u=1-\xi\)，振幅为
\(2^\delta u^\delta(1-\delta u/2+O(u^2))\)。对前两项使用端点 Fourier 公式，另一端取共轭；Gamma 递推产生显示的 \(\delta(\delta+1)/(4\pi r)\)，余项由下一阶端点估计控制。这说明这里需要的余项确实比首项低两个幂次。

对首个余弦项分部积分两次，对显示的正弦项分部积分一次，其余项按绝对值积分，得到
\[
 \int_t^\infty K_\delta(r)\,\mathrm dr
 =-\frac{a_\delta}{2\pi}t^{-\delta-1}
       \sin(2\pi t-\theta_\delta)+O(t^{-\delta-2}).
\]
所以
\[
 F_\delta(t)=1+\frac{a_\delta}{2\pi}t^{-\delta-1}
       \sin(2\pi t-\theta_\delta)+O(t^{-\delta-2})
       \quad(t\to+\infty).
\]
选择趋于无穷的一列 \(t\)，使正弦恒为一，主项最终严格大于余项，故 \(F_\delta(t)>1\)。再用 \(F_\delta(-t)=1-F_\delta(t)\)，得到负侧的严格负值。对这样一个固定 \(t\)，前面的轮廓极限还说明足够大的有限半径平均也在对应位置出现过冲。

尾积分公式使 \(|G_\delta(t)|\leq C(1+|t|)^{-\delta-1}\) 在两端成立；在有界区间它有界。因此 \(G_\delta\in L^p\) 对全部有限 \(p\geq1\) 成立，且由零点两侧趋于 \(1/2\) 和 \(-1/2\) 可知它不是零函数。将区间示性函数写成 \(H(x)-H(x-1)\)，卷积恒等式给出
\[
 B_R^\delta f(x)-f(x)
       =G_\delta(Rx)-G_\delta(Rx-R)
       \quad\text{几乎处处}.
\]
于是
\[
 R^{1/p}\|B_R^\delta f-f\|_p
       =\|G_\delta-G_\delta(\,\cdot-R)\|_p.
\]
对任意 \(G\in L^p\)，两个相距趋于无穷的平移副本满足
\(\|G-G(\,\cdot-R)\|_p\to2^{1/p}\|G\|_p\)。证明先取紧支撑截断 \(G_A=G\mathbf1_{[-A,A]}\)：当 \(R>2A\)，两支撑不交，范数的 \(p\) 次方恰为 \(2\|G_A\|_p^p\)。再由三角不等式，替换 \(G_A\) 为 \(G\) 所造成的范数误差至多为 \(2\|G-G_A\|_p\)，令 \(A\to\infty\) 即得。应用于 \(G_\delta\)，完成准确误差极限。

最后，\(B_R^\delta f\) 连续，而 \(f\) 在零点有高度一的跳跃。任意连续函数在这个跳跃两侧的本性一致误差，至少是其零点值到零和一的距离的较大者，因而不小于 \(1/2\)。这不是单个点值的误差，而是在两侧正测度邻域内都无法同时消除的误差；所以不可能有一致范数收敛。
\end{solution}
